\documentclass[twocolumn,10pt]{article}
\usepackage[utf8]{inputenc}
\usepackage[T1]{fontenc}
\usepackage{amsmath,amssymb,amsthm}
\usepackage{graphicx}
\usepackage{xcolor}
\usepackage{hyperref}
\usepackage{booktabs}
\usepackage{caption}
\usepackage{subcaption}
\usepackage{float}
\usepackage{braket}
\usepackage{multirow}
\usepackage{array}
\usepackage{mdframed}
\usepackage{tikz}
\usepackage{orcidlink}
\usepackage{enumitem}
\usetikzlibrary{arrows.meta,positioning,calc,decorations.pathreplacing,patterns,shadows.blur,fadings}
\hypersetup{
colorlinks=true,
linkcolor=blue,
citecolor=blue,
urlcolor=blue
}
\renewcommand{\hbar}{\hslash}
\newcommand{\kB}{k_{\text{B}}}
\newcommand{\Tr}{\text{Tr}}
\newcommand{\meV}{\text{~meV}}
\newcommand{\ps}{\text{~ps}}
\newcommand{\fs}{\text{~fs}}
\newcommand{\mus}{\text{~}\mu\text{s}}
\newcommand{\Hz}{\text{~Hz}}
\newcommand{\THz}{\text{~THz}}
\newcommand{\DFS}{\text{DFS}}
\newcommand{\ENAQT}{\text{ENAQT}}
\newcommand{\BQC}{B_{\text{QC}}}

\title{\textbf{The equilibrium-to-functionality gap in tubulin quantum dynamics:\\
open-system constraints, non-equilibrium mechanisms, and falsifiable tests}}
\author{
\textbf{Facundo Firmenich}$^{1,2,*}$ \orcidlink{0009-0002-6578-3811},
\textbf{Pau Firmenich}$^{1}$,
\textbf{Le\'on Firmenich}$^{1}$\\
\small $^{1}$Centro de Estudios del Sur (CEDESUR), think tank affiliated with multiple universities;
for this project, academically affiliated with Universidad Nacional Arturo Jauretche (UNAJ),
Av. Calchaquí 6200, Florencio Varela, Provincia de Buenos Aires, 1888, Argentina\\
\small $^{2}$Universitat de Barcelona (UB), Gran Via de les Corts Catalanes, 585,
08007 Barcelona, Spain\\
\small $^{*}$Correspondence: \href{mailto:f.firmenich@cedesur.org}{f.firmenich@cedesur.org}
}
\date{\today}
\begin{document}
\maketitle

% ============================================================
% GRAPHICAL ABSTRACT (TikZ)
% ============================================================
\begin{figure*}[t!]
\centering
\begin{tikzpicture}[
scale=0.95,
font=\sffamily\small,
box/.style={rounded corners=2pt, draw=black, fill=white, inner sep=4pt, align=center},
arrow/.style={->, >=Stealth, line width=0.8pt},
bridge/.style={decorate, decoration={brace, amplitude=6pt, mirror}, thick}
]
\fill[white] (-0.5,-2.5) rectangle (14.5,5.5);
\node[font=\bfseries\large, align=center] at (7,5) {Equilibrium-to-Functionality Gap:\\
Quantum Coherence in Tubulin};
\draw[->, thick] (0.5,0) -- (13.5,0) node[right] {$\log_{10}(\tau/\text{s})$};
\foreach \x/\lab/\sub in {0/-15/fs, 2/-12/ps, 4/-9/ns, 6/-6/\mus, 8/-3/ms, 10/0/s}{
\draw (\x+0.5,0.1) -- (\x+0.5,-0.1) node[below] {\lab};
\node[gray, font=\tiny] at (\x+0.5,-0.4) {\sub};
}
\fill[red!15] (0.8,-1.8) rectangle (2.8,-0.6);
\draw[red!70!black, thick] (0.8,-1.8) rectangle (2.8,-0.6);
\node[red!70!black, font=\footnotesize\bfseries, align=center] at (1.8,-1.2)
{Equilibrium\\$T_2^*\!\approx\!40$\,fs};
\node[red!70!black, font=\tiny] at (1.8,-1.6) {vibrational dephasing};
\fill[gray!20] (8.2,-1.8) rectangle (10.2,-0.6);
\draw[gray!70!black, thick] (8.2,-1.8) rectangle (10.2,-0.6);
\node[gray!70!black, font=\footnotesize\bfseries, align=center] at (9.2,-1.2)
{Neural\\10--100\,ms};
\node[gray!70!black, font=\tiny] at (9.2,-1.6) {functional window};
\draw[<->, red!70!black, thick, dashed] (2.9,-1.2) -- (8.1,-1.2);
\node[red!70!black, font=\footnotesize\bfseries, align=center] at (5.5,-0.9)
{Gap: $\sim\!10^{11}\text{--}10^{12}\!\times$};
\node[red!70!black, font=\tiny] at (5.5,-1.5) {amplification required};
\draw[orange!70!black, thick, rounded corners=3pt] (3.2,-2.2) rectangle (5.2,-1.4);
\fill[orange!20] (3.2,-2.2) rectangle (5.2,-1.4);
\node[orange!70!black, font=\footnotesize\bfseries] at (4.2,-1.8) {Fr\"ohlich};
\node[orange!70!black, font=\tiny] at (4.2,-2.0) {$\tau\!\sim\!10$--$100$\,ps};
\node[orange!70!black, font=\tiny] at (4.2,-2.15) {length-gated};
\draw[green!60!black, thick, rounded corners=3pt] (5.5,-2.2) rectangle (7.5,-1.4);
\fill[green!15] (5.5,-2.2) rectangle (7.5,-1.4);
\node[green!60!black, font=\footnotesize\bfseries] at (6.5,-1.8) {QED Cavity};
\node[green!60!black, font=\tiny] at (6.5,-2.0) {$\tau\!\sim\!0.1$--$1$\,\mus};
\node[green!60!black, font=\tiny] at (6.5,-2.15) {ordered water};
\draw[violet!70!black, thick, rounded corners=3pt] (7.8,-2.2) rectangle (9.8,-1.4);
\fill[violet!20] (7.8,-2.2) rectangle (9.8,-1.4);
\node[violet!70!black, font=\footnotesize\bfseries] at (8.8,-1.8) {Subradiance};
\node[violet!70!black, font=\tiny] at (8.8,-2.0) {$\tau\!\sim\!10$\,ps--$10$\,s};
\node[violet!70!black, font=\tiny] at (8.8,-2.15) {geometric protection};
\draw[bridge, blue!70!black] (3.0,-2.8) -- (10.0,-2.8);
\node[blue!70!black, font=\footnotesize] at (6.5,-3.3)
{$\mathcal{U} = \mathcal{K}\tau_{\text{coh}}/\tau_{\text{func}} \gtrsim 1$};
\node[box, font=\tiny] at (4.2,-3.8) {Exp.~1,5:\\2D-IR + THz};
\node[box, font=\tiny] at (6.5,-3.8) {Exp.~4:\\UV/THz cavity};
\node[box, font=\tiny] at (8.8,-3.8) {Exp.~6:\\length-resolved};
\begin{scope}[shift={(0.5,3.2)}]
\node[font=\footnotesize\bfseries] at (0,0) {Legend:};
\draw[red!70!black, thick] (0,-0.4) -- (0.8,-0.4);
\node[font=\tiny, anchor=west] at (1,-0.4) {Equilibrium (insufficient)};
\draw[orange!70!black, thick] (0,-0.7) -- (0.8,-0.7);
\node[font=\tiny, anchor=west] at (1,-0.7) {Non-equilibrium candidates};
\draw[violet!70!black, thick] (0,-1.0) -- (0.8,-1.0);
\node[font=\tiny, anchor=west] at (1,-1.0) {Experimentally observed};
\draw[gray!70!black, thick] (0,-1.3) -- (0.8,-1.3);
\node[font=\tiny, anchor=west] at (1,-1.3) {Functional target};
\end{scope}
\node[box, font=\footnotesize, text width=12.5cm, align=justify] at (7,-4.8)
{\textbf{Graphical Abstract.} The equilibrium-to-functionality gap separates
femtosecond molecular coherence (red) from millisecond neural dynamics (gray).
Three non-equilibrium mechanisms---Fr\"ohlich pumping (orange), ordered-water
QED cavity (green), and geometric subradiance (violet)---are candidate bridges.
Only subradiance is experimentally observed in related architectures
(Babcock et al. 2024). The dimensionless utility $\mathcal{U}$ quantifies
the amplification requirement ($\mathcal{K}$: effective cascade gain); 
$\mathcal{U}\!\gtrsim\!1$ defines the falsifiable experimental targets (Experiments 1, 4--6).};
\end{tikzpicture}
\caption*{\textbf{Graphical Abstract.} Visual synthesis of the equilibrium-to-functionality gap and candidate non-equilibrium bridges. See main text for quantitative details.}
\label{fig:graphical_abstract}
\end{figure*}

% ============================================================
% ABSTRACT
% ============================================================
\begin{abstract}
Biological quantum effects are no longer exceptional in principle: room-temperature
photosynthetic complexes display ultrafast electronic coherence and longer-lived
vibrational coherences, avian magnetoreceptors exploit radical-pair spin dynamics,
and enzyme catalysis can involve tunnelling. The unresolved biological-physics
question is sharper: when can short molecular quantum dynamics cross the gap to
slower cellular or neural function? We formalize this mismatch as an
\emph{equilibrium-to-functionality gap} and analyze tubulin and microtubules as a
stringent test case. Using a hierarchical open-quantum-system framework
(secular Lindblad limits, Redfield theory, and HEOM validation ledgers), we model
decoherence from thermal collisions, ionic perturbations, protein vibrational
modes, and structured environmental couplings. The calibrated equilibrium Ohmic
baseline gives $T_2^*\approx 39\fs$ at 310~K, with structured-bath variants remaining many orders of magnitude below neural timescales, and all equilibrium scenarios remain
many orders of magnitude shorter than neural synchronization scales
(10--100~ms). We therefore introduce a dimensionless quantum-coherence utility
criterion $\mathcal{U}(\tau_{\text{coh}},\tau_{\text{func}},\mathcal{K})$ that
separates equilibrium insufficiency from non-equilibrium possibility. Within this
framework, equilibrium protection and naive cooperative scaling fail, whereas
Fr\"ohlich pumping, ordered-water QED cavity models, and geometrically protected
subradiant states remain candidate non-equilibrium model classes with distinct
experimental signatures. The manuscript does not assume that tubulin quantum
dynamics are functionally relevant; it derives the quantitative conditions under
which they could become relevant and specifies how each candidate mechanism can
be falsified. We further compare these mechanisms against classical neural
synchronization channels and show that known neural-scale synchronization does
not require a quantum explanation. Six discriminating experimental programmes
are proposed. The framework is delivered with fully reproducible computational
artefacts, calibration ledgers, and open-source code.
\vspace{0.2cm}
\noindent\textbf{Keywords:} biological physics, quantum coherence, microtubules,
tubulin, open quantum systems, HEOM, Redfield theory, Fr\"ohlich condensation,
QED cavity, subradiance, quantum discord, neural synchronization, quantum biology
\end{abstract}

% ============================================================
% FUNDING
% ============================================================
\section*{Funding}
This research received no specific grant from any funding agency in the public, commercial, or not-for-profit sectors.

% ============================================================
% INTRODUCTION
% ============================================================
\section{Introduction: biological quantum relevance as an equilibrium-to-functionality problem}
The central question addressed here is not whether quantum mechanics is present
in biomolecules---it necessarily is---but whether molecular quantum dynamics can
remain operationally relevant across the enormous timescale separation between
femtosecond/picosecond molecular motion and microsecond-to-millisecond biological
function. We treat this separation as an \emph{equilibrium-to-functionality gap}.
The gap is especially severe for tubulin: microtubules are structurally ordered,
biologically central, and optically rich, yet they operate in warm, wet,
strongly fluctuating cellular environments. They therefore provide a stringent
rather than permissive test case for biological quantum relevance.
Our objective is deliberately conservative. We do not assume that tubulin quantum
coherence contributes to neural computation, nor do we attempt to explain
consciousness. Instead, we ask what physical conditions would be required for
molecular quantum dynamics in tubulin to become functionally consequential, and
which experiments could rule in or rule out the relevant non-equilibrium model
classes.
\subsection{Quantum coherence in biology: revised understanding (2017--2025)}
Quantum mechanics was traditionally considered irrelevant for biology beyond
bond formation. This shifted with demonstrations of quantum effects in three
established systems plus one emerging one.
\textbf{Photosynthetic light harvesting (FMO complex).} Initial reports
\cite{Engel2007,Panitchayangkoon2010} of $\tau \sim 300\text{--}600\fs$
electronic coherence at 277~K have been substantially revised. Thyrhaug
et al.\ \cite{Thyrhaug2018} demonstrated using polarization-controlled 2D
spectroscopy that long-lived quantum beats are exclusively vibrational in
origin, whereas electronic coherences dephase within 240~fs even at 77~K.
At physiological temperature, Duan et al.\ \cite{Duan2017} found electronic
coherence $\tau_{\text{elec}} < 60\fs$ at 295~K. Computational studies
\cite{Tiwari2013,Christensson2012,Tempelaar2014} showed vibrational nuclear
motion produces spectral signatures previously attributed to electronic
coherence. Thyrhaug et al.\ \cite{Thyrhaug2018,Thyrhaug2018b} confirmed
electronic coherences are ultrafast ($\sim 60\fs$) at room temperature for
high-energy excitons, extending to $\sim 500\fs$ only at cryogenic
temperatures.
\textbf{Functional implications.} Despite ultrafast electronic decoherence,
energy transfer in FMO occurs over $\sim 1$~ps timescales, suggesting
vibrational coherences and environment-assisted quantum transport
(\ENAQT) \cite{Plenio2008} contribute to efficiency even without persistent
electronic coherence \cite{Cao2020}. The quantum advantage appears to lie
in vibronic coupling \cite{Chin2013} and correlated environmental
fluctuations rather than long-lived electronic superpositions.
\textbf{Avian magnetoreception.} Radical pair spin coherence in cryptochrome
proteins exhibits $\tau \sim 1\text{--}10~\mu\text{s}$ at 300~K with
experimental validation via magnetic isotope effects
\cite{Ritz2000,Maeda2012,Gauger2011}.
\textbf{Enzymatic tunneling.} Kinetic isotope effects (KIE $\sim 3\text{--}80$)
in hydrogen transfer enzymes indicate quantum tunneling contributions to
catalysis \cite{Scrutton1999,Klinman2013}.
\textbf{Microtubule tryptophan networks.} Babcock et al.\ \cite{Babcock2024}
experimentally demonstrated UV superradiance in mega-networks ($>10^5$
tryptophan dipoles) of microtubule architectures at room temperature,
building on theoretical predictions of superradiant transport in tubulin
\cite{Celardo2019}. Superradiant states ($\tau \sim 100\fs$) coexisted with
subradiant states ($\tau \sim 10$~s), representing the first direct evidence
of collective quantum optical effects in MTs.
\subsection{The microtubule hypothesis: QED cavity and non-equilibrium models}
Microtubules are hollow cylinders (outer diameter 25~nm, inner lumen 15~nm,
length 1--50~$\mu$m) assembled from $\alpha\beta$-tubulin dimers in 13
protofilaments. Each neuron contains $\sim 10^7$ dimers.
\textbf{QED cavity model.} Mavromatos, Mershin, and Nanopoulos
\cite{Mavromatos2025,Mavromatos1997,Mavromatos2002} developed a model
treating MT interiors (containing ordered water
\cite{DelGiudice1988,Sahu2013}) as high-Q electromagnetic cavities. Strong
dipole-dipole coupling between tubulin dimers ($\mu \sim 1700$~Debye
\cite{Tuszynski1995}) and ordered water dipoles in a $\sim 1$~nm hydrophobic
wall layer provides environmental shielding. Decoherence arises primarily
from leakage of water dipole quanta, yielding:
\begin{equation}
\tau_{\text{QED-cavity}} \sim 10^{-7}\text{--}10^{-6}\text{~s}
\label{eq:cavity_tau}
\end{equation}
with scaling $\tau \propto \varepsilon^2$, where $\varepsilon \sim 80$ is
the ordered-water dielectric constant.
\textbf{Non-equilibrium (Fr\"ohlich) models.} Fr\"ohlich
\cite{Frohlich1968,Frohlich1977} demonstrated that dipole systems under
continuous metabolic pumping can form Bose-like condensates at high
temperature, explicitly violating equilibrium thermodynamics. For MT dimers
continuously supplied with GTP hydrolysis energy ($\Delta E \sim 0.42$~eV
$\approx 16\kB T$ at 310~K), coherence of a collective mode at frequency
$\omega_F$ could be sustained with:
\begin{equation}
\tau_{\text{Fr\"{o}hlich}} \sim Q / \omega_F \sim
10^{-11}\text{--}10^{-10}\text{~s}
\end{equation}
where $Q$ is the quality factor of the condensed mode.
\textbf{Controversy.} Reimers et al.\ \cite{Reimers2009} analyzed
one-dimensional polymers and concluded that Fr\"ohlich condensates are
inaccessible in a biological environment. However, Pokorn\'y et al.\
\cite{Pokorn2011} noted this analysis uses linearized equations neglecting
non-linear elastic-electric coupling central to Fr\"ohlich's original
framework, and models 1D chains rather than cylindrical MT geometry with
ordered water \cite{Mavromatos2025,Salari2011}. The discrepancy between
these analyses defines the central unresolved theoretical question.
\textbf{Recent experimental support.} Khan et al.\ \cite{Khan2024} showed
MT-stabilizing drugs delayed anesthetic-induced unconsciousness in rats
(Cohen's $d = 1.9$). Kalra et al.\ \cite{Kalra2023} demonstrated anesthetics
dampen quantum optical effects in MTs. Wiest \cite{Wiest2025} reviewed
convergent evidence that volatile anesthetics act on MT quantum states.
Saxena et al.\ \cite{Saxena2020} and Singh et al.\ \cite{Singh2021}
observed MT resonances spanning multiple neurons.
\subsection{Equilibrium vs non-equilibrium: the critical distinction}
Tegmark's critique \cite{Tegmark2000} assumed thermal equilibrium and large
spatial superpositions. Hagan et al.\ \cite{Hagan2002} corrected the spatial
assumption. The fundamental issue remains: equilibrium assumptions are
inappropriate for living systems continuously dissipating metabolic energy.
Our approach: \S\S 2--3 establish rigorous equilibrium constraints,
from the calibrated femtosecond pure-dephasing baseline to ps-scale
structured-bath benchmarks. \S 3.4 introduces a novel coherence utility framework
quantifying the amplification gap. \S 4 analyzes non-equilibrium mechanisms
(Fr\"ohlich, cavity, subradiance) as potential routes to ps--s effective
lifetimes, depending on the mechanism.
\S 5 provides a quantitative classical comparison with information-theoretic
distinguishability criteria. \S\ref{sec:experimental_discrimination} proposes six falsifiable experiments.
\subsection{The equilibrium-to-functionality gap}
Neural timescales (10--100~ms) compared with the calibrated equilibrium
baseline ($T_2^*\approx39\fs$) imply an amplification factor
$10^{11}\text{--}10^{12}$. Candidate non-equilibrium mechanisms reduce
this gap to different degrees: Fr\"ohlich pumping remains in the ps regime,
ordered-water cavity models reach the sub-\(\mu\)s--\(\mu\)s regime, and
reported subradiant tails in related architectures can extend much further. For functional relevance, one of the
following must hold: (1) extreme equilibrium protection yielding
$\tau \sim$ ms (unknown mechanism); (2) cooperative amplification by factor
$\sim N$ or $\sqrt{N}$; (3) non-equilibrium pumping sustaining coherence
(Fr\"ohlich); (4) cavity shielding by ordered water (Mavromatos et al.);
(5) cascading quantum$\to$classical amplification; (6) geometric
decoherence-free subspaces (subradiant states); (7) classical mechanisms
suffice (null hypothesis). We demonstrate equilibrium mechanisms (1--2)
fail; non-equilibrium mechanisms (3, 4, 6) are testable hypotheses.

% ============================================================
% FIGURE 1: TIMESCALE LANDSCAPE
% ============================================================
\begin{figure*}[t]
\centering
\includegraphics[width=0.95\textwidth]{figures_final/fig1_landscape.pdf}
\caption{\textbf{Landscape of biological quantum timescales and the equilibrium-to-functionality gap.}
Characteristic timescales for established biological quantum effects, the
calibrated tubulin equilibrium baseline, candidate non-equilibrium microtubule
mechanisms, and neural synchronization. The landscape separates the
femtosecond equilibrium regime from the millisecond functional regime and
identifies Fr\"ohlich pumping, ordered-water QED cavity protection, and
geometric subradiance as experimentally testable non-equilibrium bridges rather
than established functional mechanisms. The horizontal separation defines the
amplification gap quantified by
\(\mathcal{U}_{\rm phys}=\mathcal{K}\tau_{\rm coh}/\tau_{\rm func}\).}
\label{fig:landscape}
\end{figure*}

\subsection{Scope and contribution}
The contribution of the paper is a PRX Life-oriented biological-physics
framework: it turns a historically speculative topic into a constrained set of
open-system calculations, amplification bounds, non-equilibrium hypotheses, and
falsifiable experimental discriminants.
\textbf{Objective 1 (\S 2):} Establish equilibrium decoherence constraints for
tubulin using calibrated open-system models, spectral-density uncertainty, and
Redfield/HEOM validation ledgers. The target is not to prove functional quantum
relevance, but to quantify the equilibrium insufficiency baseline.
\textbf{Objective 2 (\S 3):} Derive cooperative amplification bounds, introduce
a quantum-coherence utility criterion, and analyze quantum discord in coupled
tubulin dimers. This section defines the equilibrium-to-functionality gap in a
form that can be applied beyond tubulin.
\textbf{Objective 3 (\S 4):} Treat Fr\"ohlich pumping, ordered-water cavity QED,
and geometric subradiance as non-equilibrium \emph{model classes}. They are not
assumed to be biologically realized; they are retained only insofar as they imply
quantitative, mechanism-specific signatures.
\textbf{Objective 4 (\S 5):} Compare candidate quantum mechanisms against known
classical synchronization mechanisms and make the null hypothesis explicit:
current neural synchronization data do not require a quantum explanation.
\textbf{Objective 5 (\S\ref{sec:experimental_discrimination}):} Provide six experimental programmes with explicit
positive criteria, null criteria, confounds, and statistical decision layers.
The six experiments collectively discriminate between equilibrium vibronic
(small-polaron), Fr\"ohlich condensate, ordered-water QED cavity, and
geometric subradiance scenarios, with Experiment~6 providing the first
quantitative map of the $(\tau_A, p)$ subradiance parameter space.

% ============================================================
% SECTION 2: EQUILIBRIUM DECOHERENCE
% ============================================================
\section{Equilibrium decoherence analysis}
\subsection{Two-level system model and the choice of open-system formalism}
We model tubulin dimers as quantum two-state systems (``straight'' $\ket{0}$
vs ``curved'' $\ket{1}$ conformations):
\begin{align}
\hat{H} &= \hat{H}_{\text{sys}} + \hat{H}_{\text{bath}} + \hat{H}_{\text{int}} \label{eq:H_total} \\
\hat{H}_{\text{sys}} &= \frac{\Delta E}{2}\sigma_z - \Delta \sigma_x \\
\hat{H}_{\text{bath}} &= \sum_k \hbar\omega_k (b_k^\dagger b_k + \tfrac{1}{2}) \\
\hat{H}_{\text{int}} &= \sigma_z \otimes \sum_k g_k (b_k^\dagger + b_k)
\end{align}
We adopt an Ohmic spectral density as the primary model:
\begin{equation}
J_{\text{Ohm}}(\omega) = \eta \omega e^{-\omega/\omega_c}
\label{eq:spectral_ohmic}
\end{equation}
with dimensionless coupling $\eta$ and cutoff $\omega_c$.
\textbf{Choice of master equation.}
At $T=310$~K with $\eta\sim 0.3$ and $\omega_c\sim 4.5\THz$ the relevant
dimensionless parameters are $\hbar\omega_c/\kB T \approx 0.68$ and
$2\pi\eta\kB T/\hbar\omega_c \approx 1.4$. The system therefore sits in the
\emph{intermediate-coupling / non-Markovian} regime where neither the
Born--Markov secular Lindblad equation nor the Redfield equation are
a priori exact \cite{BreuerPetruccione2002,Weiss2008}. We use three levels
of description, benchmarked against each other:
\begin{enumerate}
\item \emph{Secular Lindblad} --- analytically tractable; valid when
$|\omega_{ij}-\omega_{kl}|\gg \Gamma$. Used for closed-form scalings.
\item \emph{Time-convolutionless Redfield} (second order) --- retains
non-secular coherence transfer; used for rate tables (Table~1).
\item \emph{Hierarchical Equations of Motion} (HEOM)
\cite{Ishizaki2009,Tanimura2020} --- numerically exact for Drude--Lorentz
baths; used as benchmark in \S 2.2.5.
\end{enumerate}
All three agree to within $25\%$ on $\tau_c$ for the parameter window
considered, confirming that the order-of-magnitude conclusions of this
section are not an artefact of the master-equation choice. The largest
discrepancies arise for Ohmic baths at $\eta\gtrsim 0.5$, where HEOM
predicts a $\sim 40\%$ shorter $\tau_c$ than secular Lindblad due to
non-Markovian memory effects that transiently reduce coherence before
long-time thermalisation.

\subsubsection{Spectral density calibration and uncertainty}
The parameters $\eta$ and $\omega_c$ are drawn from protein electron
transfer studies \cite{May2011,Rego2003}: $\eta \in [0.1, 1.0]$ with central
value 0.3, and $\omega_c \in [100, 250]\text{~cm}^{-1}$ ($3\text{--}7.5\THz$)
with central value $150\text{~cm}^{-1}$ ($4.5\THz$). Xiang et al.\
\cite{Xiang2019} used comparable values for tubulin dephasing. We note these
are \textit{generic protein} parameters; tubulin-specific values should be
extracted from molecular dynamics simulations of the $\alpha\beta$-dimer in
explicit solvent \cite{Sept2003}, which we identify as priority future work.
To quantify parametric uncertainty, we compute $\tau_c$ across the full
$(\eta, \omega_c)$ parameter space and verify that the equilibrium scale
remains ultrafast. In the current public release, using the numerical code
(\S\ref{sec:code}), we performed a multi-temperature inversion to break the
degeneracy between $\eta$ and $\omega_c$. Synthetic benchmark tests confirm
that the inversion engine recovers the true parameters with absolute error
$\varepsilon_\eta = 1.59\times10^{-12}$ and relative error
$\varepsilon_{\omega_c}/\omega_c \approx 3.9\times10^{-10}$ under ideal
conditions, validating identifiability and calibration stability. The calibrated
equilibrium $T_2^*$ is consistently $39.1\fs$ at 310~K with $\eta=0.1$, adopted
as the baseline.

\textbf{Sobol sensitivity analysis.} A global variance-based sensitivity analysis (Sobol indices, $N=2000$) reveals that $\eta$ dominates the uncertainty in $T_2^*$: $S_1(\eta)=0.982 \pm 0.011$, $S_T(\eta)=0.991 \pm 0.006$. The cutoff frequency contributes marginally ($S_1(\omega_c)=0.014$), while temperature and protection factors are negligible ($<0.005$). This confirms that the equilibrium bound is structurally controlled by the low-frequency spectral density, not by high-frequency cutoffs or external shielding.

\textbf{Structured-bath cross-check.} To verify robustness beyond the Ohmic approximation, we evaluated a Drude--Lorentz spectral density $J_{\text{DL}}(\omega) = \frac{\eta_{\text{DL}} \omega_c \omega}{\omega^2+\omega_c^2}$ with $\omega_c \sim 50$ THz (typical protein Drude frequency) and $\eta_{\text{DL}}$ matched to the low-frequency dissipation of the Ohmic case. The extracted coherence time shifts to $T_2^* = 42 \pm 3$ fs, confirming that the equilibrium result is robust to spectral shape variations within a factor of two. The $42\pm3$ fs value refers to a low-frequency-matched Drude--Lorentz pure-dephasing surrogate used to test the robustness of the calibrated Ohmic $T_2^*$ baseline. It should not be confused with the structured Drude--Lorentz HEOM excitonic benchmark below, which probes population/coherence decay in a different effective model and yields ps-scale relaxation while still remaining many orders of magnitude below neural timescales.

\subsection{Decoherence sources}
\subsubsection{Thermal water collisions}
\begin{equation}
\Gamma_{\text{water}} = n \sigma v_{\text{th}} f_{\text{prot}}
\end{equation}
where $n \approx 3.3 \times 10^{28}$~m$^{-3}$ is water number density,
$\sigma \sim 10^{-19}$~m$^2$ is the effective collision cross-section,
$v_{\text{th}} \approx 640$~m/s at 310~K, and $f_{\text{prot}}$ is a
protection factor. Unprotected: $\Gamma \sim 10^{13}\Hz$. With hydrophobic
cavity and ordered water shielding ($f_{\text{prot}} \sim 10^{-3}$):
$\Gamma \sim 10^{10}\Hz$ \cite{Eaves2005}.

\subsubsection{Ionic fluctuations}
Debye-screened electrostatic noise from physiological ions
($I \approx 150$~mM):
\begin{equation}
\Gamma_{\text{ions}} \sim \frac{\mu^2 E_{\text{rms}}^2}{\hbar^2}
\tau_{\text{corr}} \sim 10^{9}\text{--}10^{10}\Hz
\end{equation}
where $\mu$ is the dipole moment, $E_{\text{rms}}$ the ionic field at the
Debye length $\lambda_D \approx 0.8$~nm, and
$\tau_{\text{corr}} \sim \lambda_D / v_{\text{ion}}$ the correlation time
\cite{Hagan2002,Xiang2019,Israelachvili2011}. Recent computational
refinements incorporating a screening factor $\delta = 0.005$ due to the
MT wall reduce the effective ionic rate to $\sim 9.3$~GHz, in agreement
with numerical simulations (\S\ref{sec:code}).

\subsubsection{Vibrational dephasing (dominant)}
Tubulin has $\sim 3 \times 10^4$ vibrational modes, many thermally excited
at 310~K:
\begin{equation}
\Gamma_{\text{vib}} = \frac{\pi}{\hbar} \int_0^\infty d\omega \, J(\omega)
\coth\!\left(\frac{\hbar\omega}{2\kB T}\right)
\label{eq:Gamma_vib_full}
\end{equation}
In the high-temperature regime ($\hbar\omega_c / \kB T \approx 0.7$ at
310~K), the analytic approximation for Ohmic $J(\omega)$ gives:
\begin{equation}
\Gamma_{\text{vib}} \approx \frac{2\pi \eta \kB T}{\hbar}
\end{equation}
For the calibrated release baseline $T = 310$~K, $\eta = 0.1$:
\begin{equation}
\boxed{\Gamma_{\text{vib}} \approx 2.6\times 10^{13}\Hz
\quad \Rightarrow \quad T_2^* \approx 39\fs}
\label{eq:tau_vib_boxed}
\end{equation}
\paragraph{Vibrational dominance hierarchy.}
\label{prop:vib_dominance}
For tubulin dimers with intrinsic bath coupling $\eta \geq 0.1$,
vibrational dephasing dominates all other equilibrium mechanisms by at
least one order of magnitude:
\begin{equation}
\Gamma_{\text{vib}} \gg \Gamma_{\text{ions}} \gtrsim
\Gamma_{\text{water}}^{(\text{prot})} \gg \Gamma_{\text{rad}}
\end{equation}
Consequently, external protection mechanisms (cavity shielding, Debye
screening) cannot improve the equilibrium coherence time beyond
$\tau \sim 10^{-14}\text{--}10^{-13}$~s in the calibrated Ohmic release,
as the dominant source is \textbf{intrinsic} to the protein.

\textit{Proof.} From the rates computed above:
$\Gamma_{\text{vib}} \sim 2.6\times 10^{13}$,
$\Gamma_{\text{water}}^{(\text{prot})} \sim 2\times 10^{9}$,
$\Gamma_{\text{ions}} \sim 9\times 10^{9}$,
$\Gamma_{\text{rad}} \sim 10^6$~Hz. Even complete elimination of water and
ionic contributions leaves
$\Gamma_{\text{total}} \geq \Gamma_{\text{vib}} \sim 10^{13}\Hz$.
\hfill $\square$

\textbf{Rigidity comparison.} Tubulin C-terminus B-factors
($\sim 50\text{--}80$~\AA$^2$ \cite{Nogales1998}) exceed FMO rigidity
($\sim 10\text{--}20$~\AA$^2$ \cite{Adolphs2006}), suggesting weaker
vibrational protection. However, crystallographic B-factors may not reflect
solution dynamics; NMR relaxation data for tubulin are needed.

\subsubsection{Comparison with photosynthesis (updated)}
Thyrhaug et al.\ \cite{Thyrhaug2018,Thyrhaug2018b} used polarization-controlled 2D
electronic spectroscopy to demonstrate that oscillatory signals previously
interpreted as electronic coherence are vibrational beatings, with
electronic dephasing completing within 240~fs at 77~K. Duan et al.\
\cite{Duan2017} reported $\tau_{\text{elec}} < 60\fs$ at 295~K. The revised
understanding:
\begin{equation}
\tau_{\text{elec}}^{\text{FMO}} \sim 60\fs, \quad
\tau_{\text{vib}}^{\text{FMO}} \sim 1\text{--}10\ps
\end{equation}
Our current calibrated Ohmic release value $T_2^* \approx 39\fs$ no longer
overlaps the FMO vibronic window directly. Any stronger FMO-style vibronic
analogy for tubulin therefore requires either a more structured bath model,
correlated environmental effects, or an explicitly non-equilibrium mechanism.
We retain the Holstein construction of \S\ref{sec:holstein} as a theoretical
organiser, but not as a numerically closed claim of room-temperature
vibronic equivalence in the present release.

\subsubsection{Robustness across spectral density models (HEOM-benchmarked)}
The present release ships the calibrated Ohmic baseline together with
validation artefacts, convergence ledgers, and manuscript-level HEOM/Redfield
benchmarks. Dedicated long-window production runs for structured Drude--Lorentz
and super-Ohmic baths are staged as versioned computational extensions rather
than overclaimed here as completed biological evidence. Alternative spectral
densities remain theoretically relevant, and the benchmarks below are used to
establish the scale and robustness of the equilibrium conclusion, not to claim
that all structured-bath production calculations are already closed.

\textbf{HEOM benchmark.} We ran the Ishizaki--Fleming HEOM
\cite{Ishizaki2009} for the Drude--Lorentz case with 4 Matsubara terms and
truncation depth $K=12$ (convergence verified by doubling both). The HEOM
population dynamics yield a coherence decay $\tau^{\text{HEOM}}_c
= 1.7\ps$, to be compared with $\tau^{\text{Redfield}}_c = 2.0\ps$ and
$\tau^{\text{sec-Lindblad}}_c = 2.3\ps$. The three methods thus agree to
$\lesssim 30\%$; the secular approximation systematically overestimates
$\tau_c$ by damping coherence transfer between system eigenstates. For the
Ohmic case at $\eta=0.5$, the spread widens to $\sim 40\%$
(HEOM: $0.7\ps$; Redfield: $0.9\ps$; sec-Lindblad: $1.2\ps$), confirming
that the intermediate-coupling regime warrants exact treatment for any
precision claim below a factor of two. None of these corrections alter
the order-of-magnitude result $\tau_c\sim 1\text{--}5\ps$.
The full parameter sweep ($\eta \in [0.1,1]$, $\omega_c \in [3,7.5]\THz$)
gives $\tau \in [0.3, 20]\ps$, with the geometric mean
$\bar{\tau} \approx 3\ps$ (see Supplementary Material for contour plots;
Fig.~\ref{fig:landscape} for a visual summary across models).

\textbf{Small-$N$ HEOM/Redfield validation ledger.}
In the accompanying validation artefacts, hierarchy-depth refinement at
500~fs shows monotonic observable-level contraction across populations,
coherences, and Frobenius-norm distances:
\begin{center}
\footnotesize
\begin{tabular}{lccc}
\toprule
Comparison & $dPop$ & $dCoh$ & $dFrob$ \\
\midrule
$N_C:3\to4$, $N_k=1$ & $5.51\times10^{-2}$ & $3.57\times10^{-3}$ & $9.97\times10^{-2}$ \\
$N_C:4\to5$, $N_k=1$ & $2.90\times10^{-2}$ & $2.87\times10^{-3}$ & $4.91\times10^{-2}$ \\
$N_C:5\to6$, $N_k=1$ & $1.29\times10^{-2}$ & $1.43\times10^{-3}$ & $2.24\times10^{-2}$ \\
$N_C:6\to7$, $N_k=1$ & $5.45\times10^{-3}$ & $9.10\times10^{-4}$ & $9.45\times10^{-3}$ \\
$N_C:7\to8$, $N_k=1$ & $2.09\times10^{-3}$ & $5.06\times10^{-4}$ & $3.76\times10^{-3}$ \\
$N_k:1\to2$, $N_C=5$ & $1.24\times10^{-4}$ & $3.69\times10^{-5}$ & $1.92\times10^{-4}$ \\
\bottomrule
\end{tabular}
\end{center}
Thus, within the tested regime, $N_C$ controls the residual truncation error,
whereas $N_k$ refinement is subdominant. A full-system 1JFF calibration gives
$d_{4\to5}=9.15\times10^{-3}$, $d_{5\to6}=5.02\times10^{-3}$, and
$r(1\text{JFF})=0.549$, supporting geometric contraction while keeping the
long-window $N_C=7$ production trajectory as a staged computational closure. 
The 1JFF HEOM--Redfield discrepancy at 500~fs is much larger than the projected residual
truncation error, indicating a model-level difference between perturbative and
nonperturbative treatments in this regime. 
\textbf{Short-window extrapolative check.} The 500-fs HEOM--Redfield discrepancy is used only as a conservative indicator of model-level uncertainty, not as a substitute for a long-window production trajectory. Any extrapolation beyond 500~fs is therefore treated as a diagnostic planning calculation. It does not enter the primary equilibrium bound, which is set by the analytic FDT constraint and the explicit short-window HEOM/Redfield ledger.

% ============================================================
% FIGURE 2: HEOM BENCHMARK 1JFF
% ============================================================
\begin{figure*}[t]
\centering
\includegraphics[width=0.95\textwidth]{figures_final/heom_benchmark_1JFF.png}
\caption{\textbf{Short-window HEOM benchmark dynamics for the full 1JFF system.}
Representative HEOM trajectory used as a nonperturbative short-window
validation benchmark. Left: population redistribution from the initially
excited site over the femtosecond-to-picosecond window. Right: \(L_1\)-norm
coherence showing rapid transient build-up followed by non-exponential decay.
The figure supports the interpretation of the relevant physics as transient
open-system dynamics and illustrates why HEOM and perturbative Redfield
treatments are not quantitatively interchangeable in this parameter regime.
It is not used as a substitute for the staged long-window full-system
production trajectory.}
\label{fig:heom_1jff_benchmark}
\end{figure*}

Figure~\ref{fig:heom_1jff_benchmark} shows the corresponding short-window
HEOM dynamics for the full 1JFF system. The trajectory is used as a validation
benchmark for transient open-system behaviour, not as a long-window production
closure.

\subsection{Combined equilibrium rates}
Table~\ref{tab:decoherence_rates} summarizes all equilibrium sources.
Vibrational dephasing is intrinsic and dominant
(Proposition~\ref{prop:vib_dominance}), yielding an equilibrium coherence
time of order $10^{-14}\text{--}10^{-13}$~s in the current calibrated Ohmic
release. Using the calibrated code (\S\ref{sec:code}), the multi-temperature
inversion yields $T_2^* = 39.1\fs$ at 310~K with $\eta = 0.1$, adopted as
the baseline equilibrium coherence time for subsequent analysis.

\begin{table*}[t]
\centering
\caption{Decoherence rates in tubulin dimers at 310~K under equilibrium
assumptions. Vibrational dephasing dominates
(Proposition~\ref{prop:vib_dominance}). Non-equilibrium mechanisms analyzed
in \S 4. The \emph{Experimental status} column follows a uniform scale:
NYA = not yet attempted; AI = attempted, inconclusive; PS = partially
supported; CR = consensus-robust.}
\label{tab:decoherence_rates}
\footnotesize
\begin{tabular}{p{2.0cm}p{2.2cm}ccp{2.4cm}p{1.5cm}p{3.2cm}}
\toprule
\textbf{Source} & \textbf{Mechanism} & \textbf{Unprot.} & \textbf{Prot.}
& \textbf{Protection} & \textbf{Scale} & \textbf{Experimental status} \\
& & (s$^{-1}$) & (s$^{-1}$) & mechanism & & (refs.) \\
\midrule
Water & Collisional & $10^{13}$ & $10^{10}$ & Cavity, H$_2$O order & Ext. &
PS (Sahu 2013 \cite{Sahu2013}, ordered water) \\
\addlinespace
Ions & Electrostatic & $10^{10}$ & $10^9$ & Debye ($\lambda_D \sim 1$~nm) &
Ext. & CR (Israelachvili \cite{Israelachvili2011}) \\
\addlinespace
\textbf{Vibrations} & \textbf{Phonon coupling} & $\mathbf{10^{13}}$ &
$\mathbf{2.6\times 10^{13}}$ & \textbf{Rigidity (limited)} &
\textbf{Int.} & \textbf{AI (2D-IR proposed \S\ref{sec:exp1})} \\
\addlinespace
Radiation & Spontaneous & $10^6$ & $10^6$ & None & Negl. & CR (standard QED) \\
\midrule
\textbf{Total (eq.)} & --- & $\sim 10^{13}$ & $\sim 2.6\times 10^{13}$ & --- &
Vib-lim. & --- \\
\bottomrule
\end{tabular}
\vspace{0.2cm}
\footnotesize
\textit{Refs:} H$_2$O---\cite{Eaves2005}; Ions---\cite{Israelachvili2011,Xiang2019};
Vibrations---\cite{May2011,Rego2003,Weiss2008}; B-factors---\cite{Nogales1998}.
Robustness: the calibrated Ohmic pure-dephasing baseline remains in the
femtosecond regime, while structured-bath excitonic benchmarks remain in the
picosecond regime; all variants remain many orders of magnitude below neural
timescales. (\S 2.2.5), cross-validated against HEOM
\cite{Ishizaki2009}.
\end{table*}

% ============================================================
% SECTION 3: AMPLIFICATION PROBLEM
% ============================================================
\section{The amplification problem: equilibrium mechanisms}
\subsection{Statement and scaling analysis}
Equilibrium coherence: $T_2^* \approx 39\fs$ in the calibrated release. Neural
timescale: $\tau_{\text{neural}} \sim 10\text{--}100\text{~ms}$. Required
amplification: $F = \tau_{\text{neural}}/\tau_{\text{eq}} \sim
10^{11}\text{--}10^{12}$.

\subsection{Why $\sqrt{N}$ scaling fails}
For $N$ qubits with independent baths (GHZ-type state):
\begin{equation}
\rho_{01\cdots01}(t) \propto e^{-N\Gamma t} \quad \Rightarrow \quad
\tau_N = \frac{\tau_1}{N}
\end{equation}
Coherence \textit{worsens} with $N$ (superdecoherence).
Decoherence-free subspaces require perfect noise correlation. MT symmetry
breaking $\delta\varepsilon/\varepsilon \sim 10\text{--}50\%$
\cite{Sept2003} yields $\tau_{\DFS} \sim \hbar/\delta\varepsilon
\sim 0.1\ps$. No benefit.
Dipole--dipole quantum synchronization at 8~nm:
$K \sim \mu^2/(4\pi\varepsilon_0\varepsilon_r a^3) \sim 2\text{~GHz}$,
giving $K/\Gamma \sim 10^{-3} \ll 1$ (weak coupling, amplification $<10$).

\subsection{Dimensional equilibrium bound: FDT-rigorous form}
\textbf{Reformulation as a fluctuation--dissipation corollary.}
The original Proposition in the previous version is a scaling argument.
We now derive it as a rigorous consequence of the quantum
fluctuation--dissipation theorem (FDT; Weiss \cite{Weiss2008}, ch.~4;
Kubo \cite{Kubo1966}).
\paragraph{Pure-dephasing FDT estimate.}
\label{lem:FDT}
Let a two-level system couple to a thermal bath via
$\hat{H}_{\text{int}}=\hat\sigma_z\otimes\hat{B}$ with
$\langle\hat B\rangle=0$, the bath being in Gibbs state at temperature
$T$. Let $S_{BB}(\omega)$ be the symmetrised noise spectrum. Then the
pure-dephasing rate in the Born--Markov limit is
\begin{equation}
\Gamma_\phi = \frac{1}{\hbar^2}\, S_{BB}(0)
= \frac{1}{\hbar^2}\lim_{\omega\to 0}
\bigl[2\kB T\, \chi''(\omega)/\omega\bigr],
\label{eq:FDT_pure_dephasing}
\end{equation}
where $\chi''(\omega)$ is the imaginary part of the bath susceptibility.

\textit{Proof.} Standard: combine Redfield integral expression
\cite{BreuerPetruccione2002} with the classical form of FDT
$S_{BB}(\omega)=2\kB T\,\chi''(\omega)/\omega$, valid for
$\hbar\omega\ll\kB T$, which holds at $\omega\to 0$. \hfill$\square$
\paragraph{Equilibrium coherence ceiling.}
\label{cor:ceiling}
For any system-bath coupling with dimensionless strength
$\alpha\equiv \hbar\chi''(\omega_c)/\omega_c$ and characteristic bath
cutoff $\omega_c$, the equilibrium coherence time satisfies
\begin{equation}
\tau_{\text{coh}} \;\le\; \frac{\hbar}{\alpha\,\kB T}.
\label{eq:coh_bound}
\end{equation}

\textit{Proof.} From Eq.~(\ref{eq:FDT_pure_dephasing}), $\Gamma_\phi
\ge (\alpha\kB T)/\hbar$ whenever the spectral density is non-zero at
$\omega=0$ (Ohmic or Drude--Lorentz). $\tau_{\text{coh}}=\Gamma_\phi^{-1}$
gives the bound. Violations require either (i) gapped spectral density
($J(0)=0$, super-Ohmic $s\ge 1$), (ii) non-equilibrium drive, or (iii)
active error correction/feedback. \hfill$\square$
\paragraph{Structured-bath comparison.}
For super-Ohmic baths ($s>1$) the $\omega=0$ limit vanishes and pure
dephasing is suppressed; however, \emph{energy relaxation} (T$_1$) is then
controlled by $J(\omega_0)$ and imposes its own bound. The full \(T_2\)-relaxation ceiling at 310~K is therefore
\[
\tau\lesssim\hbar/(\alpha\kB T)\approx 25\text{--}250\fs
\quad \text{for } \alpha\in[0.1,1],
\]
with the upper end corresponding to weaker coupling, for the strict FDT pure-dephasing estimate. Structured-bath excitonic
benchmarks can display ps-scale population/coherence relaxation in a
different effective model, but this does not alter the far-below-neural
conclusion.


\subsection{Quantum coherence utility framework}
\label{sec:utility}
The utility criterion is a necessary-condition accounting device, not by
itself a complete falsification test. It quantifies the minimum product of
coherence lifetime and downstream amplification required for molecular
quantum dynamics to influence a functional process at timescale
$\tau_{\rm func}$. Because the amplification factor $\mathcal{K}$ is not yet
independently measured in MT-to-neural transduction, $\mathcal{U}_{\rm phys}$
should not be interpreted as evidence for functional quantum relevance.
Rather, it defines the required gain that any proposed mechanism must later
supply and measure independently.
\paragraph{Coherence utility.}
\label{def:utility}
\begin{equation}
\mathcal{U}_{\rm phys}(\tau_{\text{coh}}, \tau_{\text{func}}, \mathcal{K})
\equiv \frac{\mathcal{K}\,\tau_{\text{coh}}}{\tau_{\text{func}}}
\label{eq:utility_phys}
\end{equation}
where $\mathcal{K}$ is an effective amplification factor encoding
cooperative coupling ($N$ coherent units), cascading gain, and signal
transduction efficiency. Functional quantum influence requires
$\mathcal{U}_{\rm phys} \gtrsim 1$.

In what follows, we therefore distinguish
\[
\mathcal{K}_{\rm req}=\frac{\tau_{\rm func}}{\tau_{\rm coh}}
\]
from any empirical amplification factor. A mechanism is physically viable
only if an independently measured cascade gain satisfies
$\mathcal{K}_{\rm emp}\gtrsim\mathcal{K}_{\rm req}$. Until such a measurement
exists, $\mathcal{K}$ is not treated as a fitted explanatory parameter.
Mechanism-specific falsification is supplied instead by direct observables:
spectral splitting, linewidth scaling, lifetime scaling, isotope effects, and
controlled perturbation experiments.
\textbf{Equilibrium bound on $\mathcal{U}_{\rm phys}$.} From
Corollary~\ref{cor:ceiling}:
\begin{equation}
\mathcal{U}_{\text{eq}} \lesssim
\frac{\mathcal{K}\hbar}{\alpha\kB T \cdot \tau_{\text{func}}}
\end{equation}
At 310~K with $\alpha = 0.3$, $\tau_{\text{func}} = 25$~ms:
\begin{equation}
\mathcal{U}_{\text{eq}} \lesssim \mathcal{K} \times 3.3 \times 10^{-12}
\end{equation}
For $\mathcal{U}_{\text{eq}} \geq 1$: this requires
$\mathcal{K} \gtrsim 3\times10^{11}$. No known equilibrium cooperative mechanism in MTs achieves this.
Dipole coupling gives $\mathcal{K} \lesssim K/\Gamma \times N \sim
10^{-3} \times 10^4 = 10$; DFS gives $\mathcal{K} \sim 1$ (destroyed by
disorder). Therefore:
\begin{equation}
\boxed{\mathcal{U}_{\text{eq}} < 1
\quad \text{for all realistic equilibrium scenarios}}
\end{equation}
\textbf{Non-equilibrium extension.} With cavity QED protection
($\tau_{\text{coh}} \sim 0.5~\mu\text{s}$ in the calibrated release):
\begin{equation}
\mathcal{U}_{\text{cavity}} \sim
\frac{\mathcal{K} \times 5\times 10^{-7}}{25 \times 10^{-3}}
= \mathcal{K} \times 2 \times 10^{-5}
\end{equation}
Now $\mathcal{U} \geq 1$ requires $\mathcal{K} > 5 \times 10^{4}$. With
$N \sim 10^7$ dimers per neuron and cascading amplification (MT $\to$ MAP
$\to$ ion channel $\to$ membrane), $\mathcal{K} \sim 10^4\text{--}10^5$
becomes \textit{conceivable} though unproven. The utility framework thus
identifies non-equilibrium mechanisms as the critical enablers, and
$\mathcal{K}$ as the key unknown.

\textbf{Detectability extension: $\mathcal{U}_{\rm det}$.}
Physical plausibility ($\mathcal{U}_{\rm phys}$) is necessary but not sufficient. 
Experimental feasibility and inferential reliability require:
\begin{equation}
\boxed{
\mathcal{U}_{\rm det}
=
\mathcal{U}_{\rm phys}
\cdot
\beta_{\rm sel}(\text{SNR},\Delta\lambda)
\cdot
\mathcal{F}_{\rm info}(\text{SNR})
\cdot
(1-P_{\rm boundary})
\cdot
S_{\rm prep}
}
\label{eq:U_det}
\end{equation}
where $\beta_{\rm sel}\in[0,1]$ is the operational selection/detectability fraction, 
$\mathcal{F}_{\rm info} = I(\text{SNR})/I_{\max}$ is the normalized Fisher information fidelity,
$P_{\rm boundary}$ is the posterior mass concentrated at physical bounds (diagnosing prior-dominated saturation),
and $S_{\rm prep}\in[0,1]$ quantifies sample/preparation stability.
This architecture separates the physical question (``could it matter?'') from the inferential question (``can we detect and infer it without artifacts?'').

\paragraph{Spectroscopic inference quality index.}
For an experimental configuration with signal-to-noise ratio ${\rm SNR}$, critical detectability threshold ${\rm SNR}_c$, boundary posterior mass $P_{\rm boundary}$, and preparation stability $S_{\rm prep}$, we define
\[
{\rm IQI}_{\rm spec}
=
\min\left(\frac{{\rm SNR}}{{\rm SNR}_c},1\right)
(1-P_{\rm boundary})S_{\rm prep}.
\]


\textbf{Relation of $\mathcal{U}$ to existing quantum-advantage frameworks.}
The utility $\mathcal{U}$ is a \emph{generalisation}, not a reinvention.
Two comparisons situate it in the literature.
\emph{(i) Giovannetti--Lloyd--Maccone (GLM) quantum-metrology bound
\cite{Giovannetti2004}.} GLM quantify advantage as the scaling of
estimation precision with the number of probes $N$: classical shot noise
$\Delta\theta\propto N^{-1/2}$ versus Heisenberg $N^{-1}$. Our utility
reduces to the GLM form when $\tau_{\text{func}}$ is identified with the
probe-interaction time and $\mathcal{K}$ with the coherent enhancement
factor $\sqrt{N}$ or $N$:
\begin{equation}
\mathcal{U}_{\text{GLM}}\,\longleftrightarrow\,
\frac{N_{\text{coh}}\,\tau_{\text{coh}}}{T_{\text{probe}}}.
\end{equation}
$\mathcal{U}$ extends GLM by allowing $\mathcal{K}$ to encode
\emph{biological} amplification cascades (MT$\to$ion channel) that have no
counterpart in metrology.
\emph{(ii) Chin--Prior--Rosenbach--Caycedo-Soler--Huelga--Plenio
\cite{Chin2013} vibronic-ENAQT advantage.} Chin et al.\ quantify the
efficiency gain from vibrationally-assisted coherent transport as
$\eta_{\text{ENAQT}}=P_{\text{quantum}}/P_{\text{classical}}$, peaking at
intermediate bath coupling. This maps to $\mathcal{U}$ with
$\mathcal{K}=\eta_{\text{ENAQT}}$ and $\tau_{\text{func}}$ equal to the
transport time across the light-harvesting complex. For FMO at 300~K,
Chin et al.\ find $\eta_{\text{ENAQT}}\sim 2$--$3$ on a $\tau_{\text{func}}
\sim 1\ps$ timescale, giving $\mathcal{U}_{\text{FMO}}\sim 2$--$3$. This
is \emph{consistent} with the observed quantum efficiency and demonstrates
that $\mathcal{U}\gtrsim 1$ is achievable when $\tau_{\text{func}}$ is
short (picoseconds, not milliseconds). For the calibrated equilibrium tubulin baseline,
$\tau_{\text{func}}/\tau_{\text{coh}}\sim10^{11}$--$10^{12}$; even
optimistic non-equilibrium mechanisms must still overcome a severe
transduction gap. Thus the obstacle is not the mere presence of coherence,
but the mismatch between molecular and functional timescales.
Thus $\mathcal{U}$ subsumes both metrology-style ($N$-scaling) and
transport-style (ENAQT) advantages and reveals that the essential obstacle
for MTs is the timescale ratio, not an intrinsic inability to sustain
coherence.

\subsection{Quantum discord in coupled tubulin dimers}
\label{sec:discord}
The main text's claim that ``no equilibrium cooperative mechanism gives
$\mathcal{U}\ge 1$'' relies on coherence as the resource. Ollivier and
Zurek \cite{Ollivier2001} introduced \emph{quantum discord} $\mathcal{D}$,
a measure of non-classical correlations that can be non-zero even for
separable (unentangled) states and can persist beyond $\tau_c$. We
evaluate $\mathcal{D}$ analytically for $N=2$ dimers with Ohmic bath.
\textbf{Model.} Two identical two-level dimers at separation $r=8$~nm
couple via dipole--dipole interaction $J_{12}\sigma_z^{(1)}\sigma_z^{(2)}$
with $J_{12}\sim \mu^2/(4\pi\varepsilon_0\varepsilon_r r^3)\approx 8\mu$eV.
Each couples locally to an effective structured protein bath calibrated to
the optimistic ps-scale benchmark ($\tau_c\simeq1$ ps). This is deliberately
an upper-bound resource estimate. Under the stricter calibrated Ohmic
pure-dephasing baseline ($T_2^*\simeq39$ fs), the same calculation would
compress all discord lifetimes by approximately a factor of 25, giving
$\tau_{\mathcal{D}}\sim0.1$ ps rather than several ps. The discord analysis
therefore does not soften the main equilibrium no-go result; it only asks
whether non-classical correlations could persist slightly longer than
entanglement under the most favourable structured-bath benchmark considered
in this release. Starting from the maximally-discordant
state $\ket{\Psi_0}=(\ket{01}+\ket{10})/\sqrt{2}$, the reduced
two-qubit density matrix at time $t$ is, in Bell-diagonal form,
\begin{equation}
\rho(t) = \tfrac{1}{4}\bigl(\mathbb{1}+\sum_{i=xyz}
c_i(t)\,\sigma_i\otimes\sigma_i\bigr),
\end{equation}
with $c_x(t)=c_y(t)=e^{-\Gamma_\phi t}\cos(2J_{12}t/\hbar)$ and
$c_z(t)=-1$ (pure dephasing preserves $\sigma_z$ anti-correlations).
\textbf{Discord evolution.} For Bell-diagonal states,
$\mathcal{D}$ admits the closed form \cite{Luo2008}
\begin{equation}
\mathcal{D}(\rho)=\min\bigl\{\mathcal{D}_1,\mathcal{D}_2\bigr\},
\end{equation}
with $\mathcal{D}_k$ derived from the eigenvalues $\{c_i\}$. Using
$\Gamma_\phi=1/\tau_c\approx 10^{12}$~Hz:
\begin{itemize}
\item At $t=0$: $\mathcal{D}=1$ bit (maximum).
\item At $t=\tau_c\simeq 1\ps$: $c_x=c_y=e^{-1}\cos(2J_{12}\tau_c/\hbar)
\approx 0.37$, giving $\mathcal{D}\simeq 0.29$ bits.
\item At $t=10\tau_c\simeq 10\ps$: $\mathcal{D}\simeq 0.003$ bits, but
$c_z=-1$ preserves a residual classical anti-correlation.
\item At $t\to\infty$: $\mathcal{D}\to 0$ (Ohmic baths fully classicalise).
\end{itemize}
Thus, in the optimistic structured-bath benchmark, discord can persist for
$\tau_{\mathcal{D}}\sim2$--$5$ ps. In the strict Ohmic baseline, the
corresponding scale is only $\sim0.1$ ps. Neither regime approaches
functional neural timescales. The result refines the equilibrium no-go
statement: non-classical correlations may outlive entanglement by a small
factor, but they do not reopen the amplification gap.

\begin{table*}[t]
\centering
\caption{Required downstream amplification factors. These are requirements,
not assumed biological gains.}
\label{tab:K_required}
\footnotesize
\begin{tabular}{p{3.2cm}p{2.5cm}p{3.0cm}p{5.0cm}}
\toprule
\textbf{Mechanism} & $\boldsymbol{\tau_{\rm coh}}$ &
$\boldsymbol{\mathcal{K}_{\rm req}=\tau_{\rm func}/\tau_{\rm coh}}$ &
\textbf{Interpretation} \\
\midrule
Ohmic equilibrium baseline & $39$ fs & $\sim6\times10^{11}$ &
No realistic equilibrium amplification known. \\
Structured-bath benchmark & $1$--$20$ ps & $\sim10^9$--$10^{10}$ &
Still far below neural timescales. \\
Fr\"ohlich pumping & $10$--$100$ ps & $\sim10^8$--$10^9$ &
Spectroscopically interesting but functionally insufficient without extreme
measured gain. \\
QED cavity & $0.1$--$1~\mu$s & $\sim2.5\times10^4$--$2.5\times10^5$ &
Potentially compatible only if a real MT-to-channel amplification cascade is
measured. \\
Subradiance & $10$ ps--$10$ s & regime-dependent &
Potentially viable only if the long-lived component couples to functional
degrees of freedom rather than remaining optically dark. \\
\bottomrule
\end{tabular}
\end{table*}

% ============================================================
% SECTION 4: NON-EQUILIBRIUM MECHANISMS
% ============================================================
\section{Non-equilibrium mechanisms: Fr\"ohlich condensation, QED cavities,
and geometric subradiance}
\subsection{Fr\"ohlich condensation: metabolic pumping}
\subsubsection{Theoretical framework}
Fr\"ohlich \cite{Frohlich1968,Frohlich1977} showed dipole systems under
continuous energy input form coherent phonon condensates when:
\begin{equation}
\eta_{\text{pump}} \equiv
\frac{P_{\text{metabolic}}}{\kB T \Gamma_{\text{collective}}} > 1
\end{equation}
where $P_{\text{metabolic}}$ is the total power input and
$\Gamma_{\text{collective}}$ is the damping of the collective mode.
\textbf{Tubulin parameters:} Dipole $\mu \sim 1700$~Debye
\cite{Tuszynski1995}; GTP hydrolysis $\Delta E \sim 0.42$~eV at rate
$\sim 1$~s$^{-1}$/dimer; power per dimer $P \sim 6.7 \times 10^{-20}$~W.

% ============================================================
% FIGURE 3: FRÖHLICH THRESHOLD
% ============================================================
\begin{figure*}[t]
\centering
\includegraphics[width=0.95\textwidth]{figures_final/fig3_frohlich.pdf}
\caption{\textbf{Fr\"ohlich pumping threshold and collective metabolic power.}
The Fr\"ohlich mechanism requires continuous non-equilibrium pumping to exceed
the damping of a collective vibrational mode. Left: pumping parameter
\(\eta_{\rm pump}\) as a function of the effective number of coherently
participating dimers, with the condensation threshold at
\(\eta_{\rm pump}=1\). Right: comparison between total metabolic power supplied
by GTP hydrolysis and the thermal dissipation floor of a representative
collective mode. The key falsifiable prediction is length-gating: isolated
dimers remain far below threshold, whereas sufficiently long or strongly
collective microtubule segments can approach threshold only if collective
damping scales sublinearly with system size.}
\label{fig:frohlich_threshold}
\end{figure*}

\subsubsection{Condensation criterion and scaling of collective damping}
For energetic estimates, the relevant dimensionless pumping parameter is
\begin{equation}
\eta_{\rm pump}(N,\Gamma_{\rm col})
=
\frac{N P_{\rm dimer}}{\kB T\,\Gamma_{\rm col}},
\qquad
P_{\rm dimer}=r_{\rm pump}\Delta E .
\end{equation}
Using $\Delta E\simeq0.42$ eV and the illustrative upper-bound rate
$r_{\rm pump}=1\,{\rm s}^{-1}$ gives
$P_{\rm dimer}\simeq6.7\times10^{-20}$ W and, at 310 K,
\begin{equation}
\eta_{\rm pump}
\simeq
1.57\times10^{1}
\left(\frac{N}{\Gamma_{\rm col}/{\rm s}^{-1}}\right)
\left(\frac{r_{\rm pump}}{1\,{\rm s}^{-1}}\right).
\end{equation}
Thus a single dimer is far below threshold:
$\eta_{\rm pump}(N=1,\Gamma_{\rm col}=10^6\,{\rm s}^{-1})
\simeq1.6\times10^{-5}$, and for a strongly damped local mode
$\Gamma_{\rm col}\sim10^{12}\,{\rm s}^{-1}$ it is
$\sim1.6\times10^{-11}$. The threshold number of coherently pumped dimers is
\begin{equation}
N_{\rm crit}
\simeq
6.4\times10^4
\left(\frac{\Gamma_{\rm col}}{10^6\,{\rm s}^{-1}}\right)
\left(\frac{r_{\rm pump}}{1\,{\rm s}^{-1}}\right)^{-1}.
\end{equation}
Therefore Fr\"ohlich condensation is not implied by metabolic energy alone:
it requires both sufficiently large coherent participation and unusually low
collective damping.
\textbf{Critical distinction from \S 3.2.} In Fr\"ohlich condensation, $N$
dimers couple to a \textit{single collective vibrational mode} at
frequency $\omega_F$, \textit{not} to $N$ independent environmental baths.
The collective mode acts as a single effective degree of freedom with
damping $\Gamma_{\rm col}$. This avoids the superdecoherence scaling of
\S 3.2, but it does not by itself guarantee threshold crossing. For the
illustrative value $\Gamma_{\rm col}=10^6\,{\rm s}^{-1}$ and
$N\sim10^3$ dimers, one obtains
\[
\eta_{\rm pump}\simeq1.6\times10^{-2}<1.
\]
Thus micron-scale segments are below threshold unless either
$\Gamma_{\rm col}$ is much smaller, the effective pumping rate is higher,
or the coherently participating population approaches
$N_{\rm crit}\sim6.4\times10^4$.
For a 13-protofilament microtubule, this corresponds to an approximate
critical coherent length
\[
L_{\rm crit}
\simeq
\frac{8\,{\rm nm}}{13}N_{\rm crit}
\simeq
40\,\mu{\rm m}
\left(\frac{\Gamma_{\rm col}}{10^6\,{\rm s}^{-1}}\right)
\left(\frac{r_{\rm pump}}{1\,{\rm s}^{-1}}\right)^{-1}.
\]
\textbf{Sensitivity to scaling of collective damping.} The assumption that
$\Gamma_{\text{collective}}$ does not scale with $N$ is critical. If
instead the collective damping scales as $\Gamma_{\text{collective}}
= \Gamma_0 (N/N_0)^\beta$ (with $\beta \ge 0$), then the pumping
parameter becomes:
\begin{equation}
\eta_{\text{pump}} = \frac{N P}{\kB T \Gamma_0 (N/N_0)^\beta}
= \frac{P N_0^\beta}{\kB T \Gamma_0} N^{1-\beta}
\label{eq:eta_pump_beta}
\end{equation}
For $\beta = 0$, we recover the optimistic scaling $\eta \propto N$. For
$\beta = 1$, $\eta$ becomes independent of $N$; for $\beta > 1$, $\eta$
decreases with $N$.
\textbf{Perturbative estimate of $\beta$ for cylindrical MT geometry.}
We now derive a linear-response scaling estimate for $\beta$, rather than leaving it as a
free parameter.
Consider a collective vibrational mode of wavelength $\lambda_F=2\pi c_s
/\omega_F$ distributed coherently over $N$ dimers occupying a cylinder of
length $L=Na$ ($a=8$~nm the dimer spacing) and fixed radius $R=12.5$~nm.
The mode radiates into a phonon bath that, at the relevant wavelengths
$\lambda_F\gtrsim 10$~nm $\gg a$, behaves as a continuum in $d$ spatial
dimensions. Fermi's golden rule gives the decay rate into the continuum
\begin{equation}
\Gamma_{\text{col}}=\frac{2\pi}{\hbar^2}\sum_k|V_k|^2\,\delta(\omega_F
-\omega_k),
\end{equation}
with system--bath matrix element $V_k\propto \tilde\phi(k)$ where
$\tilde\phi(k)=\int d^d r\,\phi(r)e^{-i k\cdot r}$ is the form factor of
the collective mode over the $N$-dimer extent. For a normalised
box-profile $\phi(r)$:
$|\tilde\phi(k)|^2\propto N\,\mathrm{sinc}^2(kL/2)$ (along the axis) times
constant factors along the radial and azimuthal directions of the bath.
Summing with the $d$-dimensional phonon density of states
$g(\omega)\propto\omega^{d-1}$ and the selection constraint
$k\sim\omega_F/c_s$:
\begin{equation}
\Gamma_{\text{col}}(N)\propto N\,
\Bigl(\tfrac{\omega_F}{c_s}\Bigr)^{d-1}
\int dk\,\mathrm{sinc}^2(kL/2)\,\delta(\omega_F-c_s k).
\end{equation}
For $L\omega_F/c_s\gtrsim 1$ the $\mathrm{sinc}$ narrows to
$\sim 2\pi/L\propto N^{-1}$, giving the cancellation
\begin{equation}
\Gamma_{\text{col}}(N)\propto N\cdot N^{-1}\cdot\omega_F^{d-1}
=\omega_F^{d-1},
\end{equation}
i.e.\ $\beta_{\text{long-L}}=0$ (the optimistic Fr\"ohlich limit). However
this regime requires $L>\lambda_F/2$, i.e.\ MTs longer than the
mode-dependent crossover length\begin{equation}
L_*(f_F)=\frac{\pi c_s}{\omega_F}
=\frac{c_s}{2 f_F}
\simeq
10\,{\rm nm}
\left(\frac{c_s}{2\times10^3\,{\rm m/s}}\right)
\left(\frac{f_F}{10^{11}\,{\rm Hz}}\right)^{-1}.
\end{equation}
Thus, for sub-THz collective modes in the $10^{10}$--$10^{11}$ Hz range,
$L_*$ lies in the $\sim 10$--$100$ nm range, not at $10\,\mu$m. A
micron-scale crossover would correspond instead to much lower-frequency
axial modes, $f_F\sim10^8$ Hz. We therefore treat the length threshold as
mode-dependent:
\begin{equation}
L_*(f_F)=\frac{c_s}{2f_F}.
\end{equation}
Experiment~5 should consequently be interpreted as a frequency-resolved
measurement of $\beta(L,f_F)$, rather than as a fixed threshold
test.

For $L<\lambda_F/2$ the form factor is flat,
$|\tilde\phi(k)|^2\propto N^2$, and
\begin{equation}
\Gamma_{\text{col}}(N)\propto N^2\,\omega_F^{d-1}
\;\Rightarrow\;\beta_{\text{short-L}}=1.
\end{equation}
The effective bath dimensionality for a cylindrical MT embedded in
ordered water is $d=3$ for transverse radiation but reduces to
$d_{\text{eff}}\simeq 2$ for modes whose wavelength exceeds the radius
(the cylinder becomes acoustically 1D+radial). Collecting:
\begin{equation}
\boxed{\;\beta\;\in\;[0,1]\;\text{depending on}\;L/\lambda_F\;\text{and
geometry}\;}
\end{equation}
\textbf{Consequence for feasibility.} Substituting into
Eq.~(\ref{eq:eta_pump_beta}):
\begin{itemize}
\item For modes with $L\ll L_*(f_F)$, $\eta_{\rm pump}$ is effectively
$N$-independent and condensation is strongly disfavoured.

\item For modes with $L\gg L_*(f_F)$, the form-factor narrowing can produce
sublinear collective damping, making condensation possible only if
$N>N_{\rm crit}$ and $\Gamma_{\text{col}}$ is sufficiently small.
\end{itemize}
This is a \emph{testable} prediction: Fr\"ohlich condensation, if it
exists at all, must be length-gated. Experiment~5 (\S\ref{sec:exp5})
proposes a direct test via length-resolved THz linewidth.
\textbf{Caveat.} The perturbative treatment neglects nonlinear
elastic--electric coupling central to the full Fr\"ohlich dynamics
\cite{Pokorn2011}; nonlinearities can self-trap the collective mode and
effectively reduce $\beta$ below the linear estimate. The bound
$\beta\in[0,1]$ should therefore be read as a \emph{linear-response}
result; the true value is likely to lie near the lower end in the
condensed phase.
For the calibrated phenomenological choice $\Gamma_{\text{coll}}=10^6$~Hz,
the explicit thresholds are much larger than the heuristic values used in
earlier drafts: for $\beta = 0$, $N > N_{\text{crit}} \approx 6.4\times 10^4$;
for $\beta = 0.2$, $N \gtrsim 10^6$; for $\beta = 0.5$, $N \gtrsim 4\times 10^9$; for
$\beta = 1$, no finite $N$ suffices. The feasibility of Fr\"ohlich
condensation therefore hinges on the microscopic origin of $\beta$, which
is addressable via experiments measuring the quality factor of the
collective mode as a function of system size.
\textbf{Predicted coherence:} The condensed mode oscillates at
$\omega_F \sim 10^{10}\text{--}10^{11}\Hz$ (sub-THz), with coherence time
$\tau \sim Q/\omega_F$. For quality factor $Q \sim 1\text{--}10$:
\begin{equation}
\tau_{\text{Fr\"{o}hlich}} \sim 10^{-11}\text{--}10^{-10}\text{~s}
\quad (10\text{--}100\ps)
\end{equation}
Still $\sim 10^{8}\text{--}10^{10}$ short of ms-scale, although
$10^{2}\text{--}10^{3}$ longer than the calibrated femtosecond equilibrium
baseline.
\textbf{Interpretative consequence.}
The corrected length scale does not by itself prove that Fr\"ohlich
condensation is impossible: for sub-THz modes, most polymerized MTs satisfy
$L\gg L_*(f_F)$. Rather, it removes the earlier micron-scale length-gating
argument as a natural biological rescue mechanism. In the physically relevant
$10^{10}$--$10^{11}$ Hz range, the decisive bottleneck is no longer whether
MTs are long enough to enter the favourable form-factor regime, but whether
they can host a collective mode with both sufficiently low absolute damping
and sufficiently large coherent participation. Under conservative energetic
accounting this requires
\[
N>N_{\rm crit}\simeq6.4\times10^4
\left(\frac{\Gamma_{\text{col}}}{10^6\,{\rm s}^{-1}}\right)
\left(\frac{r_{\rm pump}}{1\,{\rm s}^{-1}}\right)^{-1},
\]
which corresponds to tens of microns even under optimistic damping and pumping
assumptions, and becomes biologically implausible for larger
$\Gamma_{\text{col}}$ or reduced coherent participation. Fr\"ohlich pumping is
therefore retained only as a sharply constrained ps-scale spectroscopic
hypothesis, not as a robust bridge to neural timescales.

\subsubsection{Critical assessment: Reimers et al.}
Reimers et al.\ \cite{Reimers2009} concluded Fr\"ohlich condensates are
inaccessible in biological environments, rendering Orch OR untenable. This
analysis has both strengths (rigorous thermodynamic framework, explicit
energy calculations, \textit{PNAS} peer review) and limitations (1D
harmonic chains rather than 2D cylindrical geometry, linearized dynamics
omitting nonlinear elastic-electric coupling \cite{Pokorn2011}, no ordered
water interactions \cite{Mavromatos2025}). Our perturbative $\beta$
analysis (\S 4.1.2) is consistent with their negative conclusion for the
short-MT limit and with Mavromatos' positive prediction only in the
long-MT limit. Experiment~4 (\S\ref{sec:exp4}) and Experiment~5
(\S\ref{sec:exp5}) together provide empirical resolution.

\subsection{QED cavity model (Mavromatos et al.)}
\subsubsection{Mechanism}
The MT interior is modeled as an electromagnetic cavity with ordered water
\cite{DelGiudice1988,Sahu2013}. Dipole-dipole coupling in a $\sim 1$~nm
layer near hydrophobic walls creates a high-Q cavity. The single-dimer
vacuum Rabi coupling is:
\begin{equation}
\lambda_0 = \frac{d_{\text{trans}} \cdot E_{\text{ow}}}{\hbar}
\label{eq:lambda0}
\end{equation}
where $d_{\text{trans}}$ is the relevant transition dipole and
$E_{\text{ow}} = \sqrt{2\pi\hbar\omega_c/(\varepsilon\varepsilon_0 V)}$ is
the ordered-water vacuum field. Two distinct estimates of $d_{\text{trans}}$
are available in the literature. For optical transitions (tryptophan at
280~nm), $d_{\text{trans}} \sim 5$~Debye; for the low-frequency
conformational mode, recent structural data \cite{Nogales1998} suggest a
significantly larger value of $d_{\text{trans}} \sim 18.7$~Debye. The
collective Rabi splitting:
\begin{equation}
\Omega_R = 2\lambda_0\sqrt{N} \approx 6\THz \approx 25\meV
\label{eq:rabi_splitting}
\end{equation}
is approximately independent of MT length (since $\lambda_0 \propto
1/\sqrt{L}$ and $\sqrt{N} \propto \sqrt{L}$ cancel). Numerical validation
using the conformational dipole ($d_{\text{trans}} = 18.7$~D) and
ordered-water parameters yields $\Delta E = 14.0\meV$
(\S\ref{sec:code}), in reasonable agreement with the nominal prediction.
The corresponding wavelength splittings at 280~nm (tryptophan absorption)
are:
\begin{align}
\Delta\lambda &= \frac{\lambda^2 \Omega_R}{c}
\approx 1.6\text{~nm} \quad (\text{nominal, } \Omega_R = 6\THz) \label{eq:delta_lambda_nom} \\
\Delta\lambda &\approx 0.89\text{~nm}
\quad (\text{numerical, } \Delta E = 14.0\meV) \label{eq:delta_lambda_num}
\end{align}
The range $0.89\text{--}1.6$~nm defines the experimental target
(see \S\ref{sec:exp4}). This corresponds to a fractional splitting
\(\Delta\lambda/{\rm FWHM}\sim0.03\)--\(0.08\), underscoring why
standard absorption spectroscopy is insufficient and why derivative UV,
THz, or nonlinear spectroscopic methods are required.
Decoherence from water dipole quanta leakage yields
$\tau_{\text{cavity}} \sim 10^{-7}\text{--}10^{-6}\text{~s}$ with scaling
$\tau \propto \varepsilon^2$ \cite{Mavromatos2025}. AC impedance
measurements on MT ensembles show relaxation $\tau \sim \mu$s
\cite{Santelices2017}, consistent with this prediction.

% ============================================================
% FIGURE 4: CAVITY SCALING
% ============================================================
\begin{figure*}[t]
\centering
\includegraphics[width=0.95\textwidth]{figures_final/fig4_scaling.pdf}
\caption{\textbf{Scaling predictions of the ordered-water QED cavity model.}
The ordered-water cavity hypothesis predicts an approximately length-stable
collective splitting over the experimentally relevant microtubule range, while
the protected lifetime depends strongly on the effective dielectric constant of
the ordered-water environment. Left: predicted splitting as a function of
microtubule length. Right: \(\tau_{\rm cavity}\propto\varepsilon^2\) scaling.
These trends define direct falsification targets: absence of the predicted
spectral splitting, or failure of the dielectric/hydration dependence under
controlled conditions, would strongly constrain the cavity model class.}
\label{fig:cavity_scaling}
\end{figure*}

\subsubsection{Critical evaluation}
\textbf{Strengths:} accounts for experimentally confirmed ordered water
\cite{Sahu2013}; predicts specific spectroscopic signatures (Rabi splitting
in the range 0.89--1.6~nm, testable via Experiment~4); explains anesthetic
effects on MT quantum states \cite{Kalra2023}; $\varepsilon^2$ scaling
provides clear parametric dependence.
\textbf{Open questions:} ordered water structure is poorly characterized
experimentally; classical MD simulations \cite{Moix2017} show no
large-scale coherent excitations (though these do not include cavity
quantum effects); both predicted splittings are below typical protein UV
absorption linewidths ($\sim 20\text{--}30$~nm FWHM), requiring
second-derivative spectroscopy or 2D electronic spectroscopy for resolution
(\S\ref{sec:exp4}).

\subsection{Geometric subradiance as a third non-equilibrium channel}
\label{sec:subradiance}
A mechanism absent in earlier treatments---but directly supported by the
Babcock~et~al.\ \cite{Babcock2024} observations---is the
\emph{geometrically-protected subradiant subspace}.
\textbf{Dicke framework.} $N$ two-level emitters coupled to a common
radiation field split into symmetric (superradiant, $\Gamma_S\sim N\Gamma_0$)
and antisymmetric (subradiant, $\Gamma_A\sim \Gamma_0/N$ or smaller)
manifolds \cite{Dicke1954,Gross1982}. In an ideal Dicke ladder, the lowest
antisymmetric state has radiative decay $\Gamma_A\to 0$ and is
\emph{decoherence-free by pure symmetry}.
\textbf{Microtubule relevance.} Babcock~et~al.\ \cite{Babcock2024} directly
observed in tryptophan mega-networks of MT architectures two coexisting
regimes: superradiant states with $\tau_S\sim 100$~fs and subradiant
states with $\tau_A\sim 10$~s---a dynamic range of $10^{14}$, among the
largest ever reported in a room-temperature biological system. This observation is not a violation of the radiative lifetime expectations
discussed above, because subradiance protects against radiative decay by
destructive interference. It does not, by itself, eliminate vibrational or
ionic pure dephasing; the effective lifetime is therefore controlled by
symmetry-breaking mixing, treated below.
\textbf{Residual dephasing.} The subradiant state is not genuinely
immortal: pure dephasing from vibrations and ions couples the
symmetric--antisymmetric manifolds at rate
$\Gamma_{\text{mix}}\sim (\delta\omega/\omega)^2\,\Gamma_\phi$, where
$\delta\omega/\omega\sim 10\text{--}50\%$ is the MT energetic disorder.
Using $\Gamma_\phi\sim 10^{12}$~Hz this gives $\Gamma_{\text{mix}}
\sim 10^{10}\text{--}10^{11}$~Hz, i.e.\ $\tau_A^{\text{effective}}
\sim 10$~ps. The 10-s lifetime in Babcock et al.\ must therefore be
explained either by (a) reduced disorder in specific tryptophan clusters,
(b) dynamical symmetrisation by the cavity QED field, or (c) coupling to
non-radiative dark states of vibrationally hot tryptophan.

\textbf{Subradiance threshold: a falsification map.}
The coherence utility criterion $\mathcal{U}_{\text{sub}}\ge 1$ requires
$p N_{\text{eff}} \tau_A \ge \tau_{\text{func}}$, which defines a
critical lifetime threshold
\begin{equation}
\tau^*(p) \equiv \frac{\tau_{\text{func}}}{p\,N_{\text{eff}}}
= \frac{2.5\times 10^{-2}}{p\times 10^7}
= \frac{2.5\times 10^{-9}}{p} \text{~s},
\label{eq:subradiance_threshold}
\end{equation}
with $\tau_{\text{func}} = 25$~ms and $N_{\text{eff}} = 10^7$.
Rather than quoting a single point estimate, we map three physically
distinct regimes:
\begin{center}
\footnotesize
\begin{tabular}{lccc}
\toprule
Regime & $\tau_A$ & $p$ required ($\mathcal{U}\ge 1$) & Status \\
\midrule
Disorder-dominated & $\sim 10\ps$ & $p > 250$ (impossible) & Ineffective \\
Partial geometric protection & $\sim 1~\mu$s & $p > 2.5\times 10^{-3}$ & Marginal, testable \\
Babcock-like & $\sim 10$~s & $p > 2.5\times 10^{-10}$ & Broadly accessible \\
\bottomrule
\end{tabular}
\end{center}
\paragraph{Subradiance threshold.}
\label{prop:subradiance_threshold}
Geometric subradiance crosses $\mathcal{U}=1$ if and only if
$\tau_A > \tau^*(p) \equiv \tau_{\mathrm{func}}/(p\,N_{\mathrm{eff}})$.
The Babcock et al.\ \cite{Babcock2024} observation ($\tau_A \sim 10$~s)
satisfies this condition for any $p > 2.5\times 10^{-10}$, a fraction
readily achieved by geometric selection. The unresolved quantity is the
reconciliation between this lifetime and the symmetry-breaking mixing rate
$\Gamma_{\mathrm{mix}}\sim 10^{10}\text{--}10^{11}$~Hz
(\S\ref{sec:subradiance}): the 10-s lifetime is not in conflict with
$\mathcal{U}\ge 1$ but does require either reduced disorder in specific
tryptophan clusters, dynamical symmetrisation by the cavity QED field, or
coupling to non-radiative dark states. This reconciliation is the primary
theoretical target of Experiment~6 (\S\ref{sec:exp6}), where
length-resolved linewidth scaling directly constrains which region of
the $(\tau_A, p)$ plane is physically realised.

Subradiance thus emerges as the \emph{only} currently \emph{observed}
mechanism capable of satisfying $\mathcal{U}\gtrsim 1$ in MT architectures,
and Babcock et al.\ shifts from an anecdotal data point to an
\emph{anchor in the $(\tau_A, p)$ parameter space} whose position
experiment must confirm or refute.
\textbf{Testability.} Babcock~et~al.\ used a specific polymerised MT
geometry; reproducing their result with varying MT lengths and
stabiliser-dependent geometries (taxol vs nocodazole) would test whether
the 10-s subradiant tail is a generic feature of MT architectures or a
geometry-specific artefact. Length-resolved linewidth scaling---the subject
of Experiment~6 (\S\ref{sec:exp6})---provides a direct constraint on
the $(\tau_A, p)$ region (Proposition~\ref{prop:subradiance_threshold})
and is the recommended next experimental step.

\subsection{Summary: non-equilibrium mechanisms}
\label{sec:mech_summary}
Table~\ref{tab:mech_comparison} synthesizes the three candidate non-equilibrium mechanisms, their scaling laws, critical thresholds, and experimental status.

\begin{table*}[t]
\centering
\caption{Unified comparison of candidate mechanisms. The listed times are not
all the same physical observable: the Ohmic baseline reports a pure-dephasing
time $T_2^*$, whereas structured-bath benchmarks report effective excitonic
population/coherence relaxation times. They are grouped only to show that all
equilibrium variants remain far below neural functional timescales.}
\label{tab:mech_comparison}
\footnotesize
\begin{tabular}{p{2.5cm} p{2.5cm} p{1.2cm} p{1.2cm} p{5.5cm}}
\toprule
Mechanism & $\tau$ range & $\mathcal{U}$ est. & $\mathcal{K}$ req. & Main confound / anchor \\
\midrule
Ohmic pure-dephasing baseline & $T_2^*\simeq39$ fs & $\ll1$ & $\sim 10^{11}$ & Intrinsic protein vibrational dephasing; strict equilibrium bound \\
Structured-bath excitonic benchmarks & $\sim0.3$--$20$ ps & $\ll1$ & $\sim 10^9$ & Effective excitonic population/coherence dynamics, not identical to pure-dephasing $T_2^*$ \\
Fr\"ohlich pumping & $10$--$100\,\mathrm{ps}$ & $\ll 1$ & $\sim 10^8$ & Metabolic pumping efficiency, low-damping modes \\
QED cavity & $0.1$--$1\,\mu\mathrm{s}$ & $0.1$--$1$ & $\sim 10^4$ & Radiative leakage, dipole strength uncertainty \\
Subradiance & $10\,\mathrm{ps}$--$10\,\mathrm{s}$ & $1$--$40$ & $1$--$10^9$ & Symmetry breaking, mixing with dark states \\
\bottomrule
\end{tabular}
\end{table*}

\subsection{Fisher detectability barrier and monotonic experimental metrics}
\label{sec:fisher_barrier_metrics}

A central difficulty in testing mesoscopic quantum-biological mechanisms is
that physical plausibility and experimental identifiability are not equivalent.
A mechanism may satisfy a formal timescale or coupling requirement while still
being statistically unrecoverable from spectra dominated by finite SNR,
instrumental resolution, boundary-constrained posteriors, or preparation
instability. We refer to this regime as a \emph{Fisher detectability barrier}:
the point at which the Fisher information carried by the measured observable
is insufficient to distinguish the candidate mechanism from the relevant null
model at the required parameter resolution.

This barrier is not a new physical interaction; it is an inference boundary.
It marks the transition between a regime where a mechanism-specific parameter
is identifiable from the data and a regime where model selection, posterior
mass, or fitted parameters become dominated by noise, instrumental priors, or
physical bounds. In the present context, the barrier is operationalized through
the detectability utility \(\mathcal{U}_{\rm det}\), the spectroscopic inference
quality index \({\rm IQI}_{\rm spec}\), and mechanism-specific monotonic
metrics.

The construction rests on standard statistical ingredients: Neyman--Pearson
model discrimination \cite{NeymanPearson1933}, Fisher-information/Cramér--Rao resolution limits \cite{Cramer1946,Rao1945},
spectroscopic SNR thresholds \cite{Kay1993}, Bayesian or information-criterion model
selection \cite{Schwarz1978}, and preparation-stability controls. Its purpose is not to replace
full mechanistic modelling, but to prevent a physically possible mechanism
from being interpreted as empirically supported when the experiment lacks the
information content required to identify it.

We therefore define mechanism-specific monotonic metrics that move predictably
with the underlying physics:

\begin{itemize}
\item \(\mathcal{M}_{\rm length}=d\log\Gamma/d\log L\), the length-linewidth
slope used to test Fr\"ohlich collective damping and the exponent \(\beta\).
\item \(\mathcal{M}_{\rm cav}=\Delta\lambda_{\rm obs}/\Delta\lambda_{\rm pred}\),
the observed-to-predicted splitting ratio for the ordered-water QED cavity
hypothesis.
\item \(\mathcal{M}_{\rm sub}=P(\tau_A>\tau^*)\) or \(\tau_A/\tau_S\), the
subradiant protection strength relative to the utility threshold.
\item \(\mathcal{M}_{\rm vib}=A_{\rm cross}/A_{\rm diag}\) as a function of
temperature and microtubule length, separating equilibrium vibronic coupling
from non-equilibrium transport enhancement.
\end{itemize}

These metrics replace binary ``detected/undetected'' claims with continuous,
mechanism-specific observables. A positive result must move the appropriate
metric in the predicted direction while surviving controls for SNR, model
selection, boundary effects, and sample stability.



% ============================================================
% SECTION 5: CLASSICAL SYNCHRONIZATION
% ============================================================
\section{Classical synchronization mechanisms}
\subsection{Inhibitory network gamma (ING/PING)}
Mathematical model \cite{Borgers2003,Brunel2003,Buzsaki2012}:
\begin{align}
\tau_m \frac{dV_i}{dt} &= -V_i + I_{\text{ext}} - g_{\text{syn}}
\sum_j s_j \\
\tau_{\text{syn}} \frac{ds_i}{dt} &= -s_i + \sum_k \delta(t - t_k^i)
\end{align}
Synchronization precision: $\sigma_{\text{ING}} \approx \sqrt{\tau_{\text{syn}}
\tau_{\text{refrac}}/N} \approx 0.14$~ms. Experimental validation: Fries
et al.\ \cite{Fries2001} (V1 gamma $\sigma \sim 3\text{--}5$~ms); Buschman
\& Miller \cite{Buschman2007} (PFC $\sigma \sim 1\text{--}3$~ms); Cardin
et al.\ \cite{Cardin2009} (optogenetic gamma $\sigma \sim 2\text{--}4$~ms).

\subsection{Gap junctions and ephaptic coupling}
Gap junctions: $\tau_{\text{gap}} \sim C_m/g_{\text{gap}} \sim 0.2$~ms,
precision $\sigma \sim 0.06$~ms \cite{Galarreta1999,Bennett2004}. Ephaptic
fields ($E \sim 1\text{--}10$~mV/mm): $\sigma \sim 1\text{--}5$~ms
\cite{Anastassiou2011,Frohlich2010}.

\subsection{Communication through coherence and dendritic computation}
Fries' communication through coherence (CTC) \cite{Fries2015} proposes
gamma synchronization gates information routing. Active dendritic
computation \cite{Poirazi2020} implements complex integrations on
millisecond scales. Together these provide a comprehensive classical
framework for neural coordination without invoking quantum effects.

\subsection{Quantitative comparison}
\begin{table}[h]
\centering
\caption{Classical mechanisms vs quantum necessity.}
\label{tab:classical_quantum}
\small
\begin{tabular}{p{2.8cm}p{1.8cm}p{2.2cm}c}
\toprule
\textbf{Phenomenon} & \textbf{Required} & \textbf{Classical} &
\textbf{Sufficient?} \\
\midrule
Gamma (40~Hz) & $\sim 10$~ms & ING/PING & \textbf{Yes} \\
Local ($<1$~mm) & $\sim 0.5$~ms & Gap junction & \textbf{Yes} \\
Long-range (cm) & $\sim 2$~ms & Ephaptic/CTC & \textbf{Yes} \\
Sub-0.1~ms (hyp.) & $< 0.1$~ms & Unknown & \textbf{Open} \\
\bottomrule
\end{tabular}
\end{table}

\subsection{Information-theoretic framework}
To quantify the informational capacity of classical vs quantum neural
channels, we adopt the mutual information framework
\cite{Borst1999,Strong1998}. Classical spike-timing precision of
$\sigma \sim 1$~ms at firing rate $r \sim 50$~Hz yields an information
rate per neuron:
\begin{equation}
R_{\text{classical}} \sim r \log_2\!\left(\frac{1}{r\sigma}\right)
\sim 50 \times \log_2(20) \approx 216 \text{~bits/s}
\end{equation}
In bursting regimes ($r \sim 10^3$~Hz), $R_{\text{classical}}$ can approach $10^4$~bits/s, narrowing the gap with the quantum transduction ceiling $B_{\text{QC}} \lesssim 10^4$~Hz set by ion channel kinetics. Consequently, any quantum advantage would not reside in raw bitrate, but in \textit{coding efficiency}, \textit{energy minimization per bit}, or \textit{sub-millisecond temporal precision} undetectable by classical Poisson-like spike trains. The effective information rate is therefore:
\begin{equation}
\Delta R_{\text{eff}} = \min\!\bigl(r_{\text{coh}},\;B_{\text{QC}}\bigr)
\cdot \log_2 d_{\text{eff}}
\;\lesssim\; 10^4 \cdot \log_2 2 \;=\; 10^4\;\text{bits/s},
\label{eq:DeltaR_eff}
\end{equation}
where $d_{\text{eff}}$ is the accessible post-transduction Hilbert space
dimension. This result holds regardless of $r_{\text{coh}}$: even
$r_{\text{coh}} = 10^{13}$~s$^{-1}$ yields $\Delta R_{\text{eff}}
\lesssim 10^4$~bits/s, comparable to the classical rate
$R_{\text{classical}} \approx 216$~bits/s above. Thus quantum coherence
offers no \emph{rate} advantage under equilibrium conditions by the
biophysical ceiling set by $B_{\text{QC}}$.
The $B_{\text{QC}}$ parameter connects the utility framework of
\S\ref{sec:utility} to the present information-theoretic analysis in a
formally precise way. From Definition~\ref{def:utility},
$\mathcal{U} \gtrsim 1$ requires $\mathcal{K} \gtrsim
\tau_{\text{func}}/\tau_{\text{coh}}$, where $\mathcal{K}$ encodes
cascading gain from the MT to the ion channel (MT~$\to$~MAP~$\to$~channel).
This cascade gain and the transduction bandwidth are related by
\begin{equation}
\mathcal{U} \gtrsim 1 \;\iff\;
B_{\text{QC}} \;\gtrsim\; \frac{1}{\mathcal{K}_0\,\tau_{\text{func}}},
\label{eq:U_BQC}
\end{equation}
where $\mathcal{K}_0$ is the MT~$\to$~ion-channel cascade factor.
Equation~(\ref{eq:U_BQC}) shows that $\mathcal{K}$ in
Definition~\ref{def:utility} and $B_{\text{QC}}$ in this section are
the same physical quantity viewed from two different angles: the
utility framework quantifies the amplification requirement;
the information-theoretic framework locates the bottleneck.
Experiment~3 (\S\ref{sec:exp3}) provides a direct upper bound on
$B_{\text{QC}}$: if tau-knockdown-mediated MT shortening alters
$\sigma_{\text{sync}}$ beyond the classical structural effect,
the difference directly constrains $\mathcal{K}_0$ and hence
$B_{\text{QC}}$ from above.
With cavity-protected coherence ($\tau \sim \mu$s), quantum states could
encode higher-dimensional information per coherent event via entanglement
across $N$ dimers \cite{Giovannetti2004}, potentially providing advantages
in coding efficiency rather than raw rate. This remains speculative and
untested.

\subsection{Distinguishability criteria}
We identify five experimental signatures that would \textit{necessitate}
quantum explanations beyond classical neural mechanisms:
(i)~sub-0.1~ms synchronization precision exceeding gap junction limits;
(ii)~violation of Bell-type inequalities in neural correlation
measurements (requiring quantum channels between neurons);
(iii)~temperature-dependent kinetic isotope effects with KIE $> 7$ and
non-Arrhenius temperature dependence in MT-dependent processes;
(iv)~MT-specific disruptions affecting cognitive functions not explainable
by structural roles alone (e.g., effects of nocodazole on synchronization
precision, not just amplitude);
(v)~isotope effects on consciousness or behavior (e.g., deuterium effects
on anesthetic potency beyond classical mass effects).
None of these have been demonstrated to date, establishing that classical
models remain sufficient for all current observations.

% ============================================================
% SECTION 7: VIBRONIC HOLSTEIN ANALOGUE
% ============================================================
\section{Vibronic Holstein analogue for tubulin}
\label{sec:holstein}
The qualitative analogy FMO$\leftrightarrow$tubulin
($\tau_v^{\text{FMO}}\leftrightarrow\tau_c^{\text{tubulin}}$) invoked in
\S 2.2.4 is strengthened by making the Hamiltonian explicit. We adopt a
minimal Holstein-type vibronic model \cite{Holstein1959,Chin2013} for the
$\alpha\beta$-tubulin dimer:
\begin{align}
\hat{H}_{\text{H}} &= \sum_{i=1,2}\varepsilon_i\,c_i^\dagger c_i
- J(c_1^\dagger c_2 + c_2^\dagger c_1) \\
&+ \sum_i \hbar\Omega_0\,b_i^\dagger b_i
+ g\sum_i c_i^\dagger c_i(b_i+b_i^\dagger)
\nonumber
\end{align}
where $c_i$ are electronic amplitudes of the conformational two-level
system on dimer $i$, $b_i$ is a local vibrational mode of frequency
$\Omega_0$, and $g$ is the electron--phonon coupling. The
Huang--Rhys factor $S=(g/\hbar\Omega_0)^2$ measures vibronic coupling
strength.
\textbf{Parameter mapping for tubulin.} From molecular-dynamics
electrostatics on the tubulin $\alpha\beta$-dimer \cite{Sept2003,Xiang2019}
and the conformational energetics of the curved--straight transition, we
estimate:
\begin{itemize}
\item $\varepsilon_2-\varepsilon_1\equiv\Delta\epsilon\approx 0.15$~eV
(GTP-state energy gap);
\item $J\approx 8\text{--}20~\mu$eV (dipolar nearest-neighbour within a
protofilament at 8-nm spacing);
\item $\hbar\Omega_0\approx 6\text{--}8$~meV (typical low-frequency
protein mode $\sim 50$--$70$~cm$^{-1}$);
\item $S\approx 0.5\text{--}1.5$ (intermediate vibronic regime, as in FMO).
\end{itemize}
With $J \approx 8\text{--}20~\mu$eV and $\hbar\Omega_0 \approx 6\text{--}8$~meV,
the ratio $J/\hbar\Omega_0 \sim 10^{-3}$ places tubulin firmly in the
\emph{small-polaron (Marcus) regime}, not the vibronic-ENAQT window.
The ENAQT window identified by Chin et al.\ \cite{Chin2013} requires
$J \sim \hbar\Omega_0$; for FMO this ratio is of order unity
($J_{\text{FMO}} \sim 100~\mu$eV, $\hbar\Omega_0^{\text{FMO}} \sim 100\text{--}200~\mu$eV).
The Holstein renormalisation gives an effective inter-dimer coupling
$J_{\text{eff}} = J\,e^{-S} \approx 0.1\text{--}0.6\,J$ for
$S \approx 0.5\text{--}1.5$, which further suppresses transfer.
This is not a failure of the analogy---it is a \textbf{quantitative discriminator}.
Two physically distinct predictions follow, and it is essential not to conflate them.
\emph{Electronic transport efficiency.} The coherent inter-dimer electronic
transfer rate is suppressed relative to the FMO ENAQT-optimal case by
$(J/\hbar\Omega_0)^2 \sim 10^{-6}$. This is a statement about
\textit{functional} quantum transport, not about spectroscopic observability.
In the Marcus regime, any net vibronic enhancement of population transfer is
negligible compared to classical thermally-activated hopping.
\emph{2D-IR cross-peak visibility.} The cross-peak amplitude in 2D-IR is
controlled independently by the Huang--Rhys factor
$S = (g/\hbar\Omega_0)^2$, not by $J$. With $S \approx 0.5$--$1.5$
and vibronic mixing angle $\theta \approx 0.3$, the cross/diagonal ratio
is $\sin^2(2\theta)\,e^{-\tau/\tau_c} \approx 3$--$5\%$ at $\tau = 3\ps$
(\S\ref{sec:exp1}, Eq.~below)---a signal that is \textit{detectable} in
principle, but that in the small-polaron regime does \textit{not} indicate
functional ENAQT transport. Observing a $3$--$5\%$ cross-peak is
consistent with intermediate vibronic coupling $S \sim 1$ and is orthogonal
to the question of whether $J/\hbar\Omega_0 \sim 1$ (ENAQT).
These two quantities are orthogonal: a system can be in the Marcus regime
($J/\hbar\Omega_0 \ll 1$, transport suppressed by $10^{-6}$) and
simultaneously produce a detectable 2D-IR cross-peak
(cross/diagonal $\sim S \sim$ few percent). The combination of both is the
equilibrium prediction for tubulin.
\emph{Non-equilibrium Fr\"ohlich pumping provides a specific resolution
pathway}: if the condensed collective mode renormalizes the effective
site coupling toward $J_{\text{eff}}^{(\text{pump})} \sim \hbar\Omega_F$,
the system crosses from the Marcus into the ENAQT window. This would
manifest as a \textit{length-dependent increase} of the cross-peak
amplitude \textit{beyond} the equilibrium $3$--$5\%$ baseline---distinguishable
from the equilibrium 3--5\% by its MT-length dependence, directly testable
by combining Experiments~1 and~5 (\S\ref{sec:exp1}, \S\ref{sec:exp5}).
Under equilibrium conditions, the intermediate $S$ value is consistent
with (not explanatory of) vibronic transport, and any FMO analogy requires
explicit non-equilibrium support.
\textbf{Prediction.} Under this mapping the Holstein model predicts four
distinguishable effects:
(i)~Equilibrium coherent inter-dimer \textit{electronic} transport is
suppressed relative to FMO by $(J/\hbar\Omega_0)^2 \sim 10^{-6}$;
structured-bath or non-equilibrium corrections are therefore required for
any vibronic transport claim beyond the calibrated Ohmic baseline.
(ii)~The 2D-IR cross-peak cross/diagonal amplitude ratio at the first beat
maximum ($\tau = \pi/\omega_{12} \approx 3\ps$) is $\approx 3$--$5\%$,
set by $S \approx 1$ and $\theta \approx 0.3$ independently of $J$.
This signal is in principle detectable but is \textit{not} diagnostic of
functional ENAQT transport: it confirms intermediate vibronic coupling,
not transport efficiency.
(iii)~If Fr\"ohlich pumping renormalizes
$J_{\text{eff}}^{(\text{pump})} \to \hbar\Omega_F$, the system crosses
from the Marcus into the ENAQT regime. This would manifest as a
\textit{length-dependent increase} of the cross-peak amplitude
\textit{above} the equilibrium $3$--$5\%$ baseline---the primary joint
discriminant of Experiments~1 and~5.
(iv)~A length-independent cross-peak at the equilibrium $3$--$5\%$ level
with no dependence on MT length is the \textit{null result} for ENAQT: it
confirms vibronic coupling (detectable) but refutes functional
non-equilibrium transport enhancement (absent).
\textbf{Interpretative shift.} The Holstein mapping reframes the
central question from ``is coherence present?'' to
``is coherence \textit{functional}?'' The equilibrium prediction is
precisely that both answers can be independently true: a
$\sim 3$--$5\%$ 2D-IR cross-peak (vibronic coupling, controlled by $S$)
coexists with a $(J/\hbar\Omega_0)^2 \sim 10^{-6}$ suppression of
ENAQT-type electronic transport (controlled by $J/\hbar\Omega_0$).
Observing the cross-peak does not establish functional transport;
\textit{not} observing a length-dependent increase in the cross-peak under
varied MT length (Experiments~1+5) would falsify the non-equilibrium
Fr\"ohlich-mediated ENAQT scenario. The functional question that remains open
is whether Fr\"ohlich or other non-equilibrium pumping drives
$J_{\text{eff}}^{(\text{pump})} \to \hbar\Omega_F$, promoting the system
into the ENAQT window---which is precisely what the joint Experiments~1
and~5 (\S\ref{sec:exp1}, \S\ref{sec:exp5}) are designed to test.
\textbf{Limitations.} The Holstein reduction retains a single vibrational
mode per dimer; a realistic treatment requires the full vibrational
density of states, which for tubulin spans $10\text{--}3000$~cm$^{-1}$. A
multimode Brownian oscillator extension \cite{Mukamel1995} or a HEOM
treatment with Matsubara-expanded correlated baths is left as priority
follow-up.

% ============================================================
% SECTION 6: EXPERIMENTAL DISCRIMINATION
% ============================================================
\section{Experimental discrimination programme}
\label{sec:experimental_discrimination}
The experimental goal is not to confirm a preferred quantum narrative, but to
separate candidate mechanisms by signatures they cannot all share. Table
\ref{tab:mechanism_discrimination} summarizes the primary positive and null
criteria used throughout the experiments below.
\begin{table*}[t]
\centering
\caption{Mechanism-specific discriminating signatures. The table is intended as
a falsification map: each candidate non-equilibrium mechanism survives only if
its own positive criterion is observed while the listed null criteria and
classical confounds are controlled.}
\label{tab:mechanism_discrimination}
\footnotesize
\begin{tabular}{p{2.6cm}p{3.0cm}p{3.5cm}p{3.0cm}p{3.2cm}}
\toprule
Mechanism & Primary signature & Positive criterion & Null criterion & Main confound \\
\midrule
Equilibrium vibronic coherence & 2D-IR cross-peaks & Temperature-dependent cross/diagonal ratio with ps-scale vibronic decay & No statistically resolved cross-peak above thermal background & Vibrational congestion and sample heterogeneity \\
Fr\"ohlich pumping & Length-dependent THz linewidth & $\beta\lesssim0.1$ with long-MT length gating & $\beta\gtrsim0.9$ or no length dependence & Bundling, damping heterogeneity, preparation variability \\
Ordered-water QED cavity & UV/THz spectral splitting & Predicted UV derivative doublet or THz cavity mode at high SNR & No feature at predicted band under controlled hydration & Turbidity, scattering, nonspecific absorption \\
Geometric subradiance & Collective lifetime/linewidth scaling & Lifetime or linewidth scaling with tryptophan-network geometry; $(\tau_A,p)$ above $\tau^*(p)$ curve & Independent-emitter scaling; $\tau_A(L)$ flat and short & Disorder and symmetry breaking \\
Classical cytoskeletal route & Tau/actin/tubulin perturbation controls & MT-specific effect not reducible to mechanics or trafficking & No specific MT contribution & Structural disruption and indirect network effects \\
\bottomrule
\end{tabular}
\end{table*}

\subsection{Validation hierarchy L1--L4}
To ensure inferential closure, all experimental proposals are structured under a four-tier validation architecture:
\begin{table*}[t]
\centering
\caption{Validation hierarchy L1--L4.}
\footnotesize
\begin{tabular}{p{1.5cm}p{3.5cm}p{8.5cm}}
\toprule
Level & Scope & KwanTube Implementation \\
\midrule
L1 & Analytic structure & FDT bounds, $\mathcal{U}_{\rm phys}$, Holstein mapping, monotonic metric definitions \\
L2 & Model calibration & SNR/Fisher calibration, synthetic spectra generation, BIC/power surfaces \\
L3 & Internal cross-check & HEOM/Redfield degradation, noise injection, spectral posterior predictive checks (PPC) \\
L4 & External validation & Babcock/Kalra/Sahu datasets, independent spectroscopic benchmarks \\
\bottomrule
\end{tabular}
\end{table*}
This hierarchy prevents circular interpretation, separates theoretical bounds from empirical detection limits, and ensures that claims of detectability are grounded in calibrated inference, not idealized assumptions.

\subsection{Experiment 1: Temperature-dependent 2D-IR}
\label{sec:exp1}
\textbf{Objective:} Detect conformational coherence; distinguish
vibrational from electronic origin.
\textbf{Method:} 2D-IR spectroscopy on amide~I band ($\sim 1650\text{~cm}^{-1}$)
of polymerized MTs.
\textbf{Key diagnostic} \cite{Christensson2012,Dean2017}: electronic
coherence produces linewidth $\Gamma(T) \propto T$; vibrational coherence
yields $\Gamma(T) \propto \coth(\hbar\omega_{\text{vib}}/2\kB T)$.
Temperature range: 280--340~K.
\textbf{Predictions.} Coherence exists: oscillatory cross-peaks with period
$T \sim h/\Delta E \sim 7\ps$, decay $\tau \sim 1\text{--}10\ps$. No
coherence: monotonic exponential decay. Controls: denatured tubulin,
monomer vs polymer.
\textbf{Quantitative cross-peak amplitude.} From response-function
formalism \cite{Mukamel1995} the rephasing 2D-IR amplitude at the
cross-peak $(\omega_i,\omega_j)$ scales as
\begin{equation}
A_{ij}^{(4)}(\tau)=|\mu_i|^2|\mu_j|^2\,
R_{ij}(\tau)\,\bigl|\langle ij|\hat\rho(\tau)|ij\rangle\bigr|,
\end{equation}
with the fourth-order response function
$R_{ij}(\tau)=e^{-\tau/\tau_c}\cos(\omega_{ij}\tau)$ in the secular
approximation. For the tubulin Holstein parameters of
\S\ref{sec:holstein} ($S\approx 1$, $\hbar\Omega_0\approx 7$~meV,
$\tau_c\approx 2\ps$, vibronic mixing angle $\theta\approx 0.3$), the
expected cross/diagonal amplitude ratio at the first beat maximum
$\tau=\pi/\omega_{12}\approx 3\ps$ is
\begin{equation}
\frac{A^{(4)}_{\text{cross}}}{A^{(4)}_{\text{diag}}}\Bigg|_{\tau=3\ps}
=\frac{|\mu_1|^2|\mu_2|^2}{|\mu_1|^4}\,\sin^2(2\theta)\,e^{-\tau/\tau_c}
\approx 0.18\,e^{-3/2}\approx 4\times 10^{-2}.
\end{equation}
Thus a cross/diagonal ratio of $\sim 3\text{--}5\%$ with a $\sim 3\ps$
oscillation period is the quantitative target. Below $\sim 1\%$ the null
hypothesis (no conformational vibronic coherence) is preferred.
Signal-to-noise requirements follow from standard pulse-shaper 2D-IR
performance \cite{Hamm2011}: $>10^4$ scans per temperature point are
needed at amide-I cross-peak level.
\textbf{Power analysis:} $n=5$ replicates/condition, Cohen's $d \sim 0.8$
detectable at $\alpha = 0.05$ with $\sim 80\%$ power.
\textbf{IQI estimate:} $\text{IQI}_{\rm spec} \approx 0.92$ (assuming $S_{\rm prep}\approx 0.95$, $\text{SNR}/\text{SNR}_c > 1$, $P_{\rm boundary} < 0.05$).

\subsection{Experiment 2: Kinetic isotope effects}
Deuterated modulators (colchicine-d$_{10}$, taxol-d$_{30}$, GTP-d$_{12}$).
Measure KIE $= k_H/k_D$ at multiple temperatures. Tunneling: KIE $> 2$,
increases with decreasing $T$. Classical: KIE $\sim 1.2\text{--}1.5$,
Arrhenius-consistent. Statistics: two-way ANOVA (isotope $\times$
temperature), $n = 3 \times 3$ replicates.
\textbf{IQI estimate:} $\text{IQI}_{\rm spec} \approx 0.88$ (standard biochemical assay precision, high $S_{\rm prep}$).

\subsection{Experiment 3: MT integrity and synchronization}
\label{sec:exp3}
MEA recordings of cortical cultures with nocodazole (MT disruptor), taxol
(stabilizer), and cytochalasin~D (actin control). Metrics: phase-locking
value (gamma band), cross-correlation timing precision $\sigma$,
spike-time tiling coefficient. Power: for 30\% PLV change (SD $\sim 15\%$),
$n = 8$ cultures at 80\% power ($\alpha = 0.05$, paired $t$-test).
\textbf{Negative control via tau knockdown.} The nocodazole--taxol--cytD
triad leaves open a confound: both nocodazole and taxol alter MT
structural abundance, which has direct non-quantum effects on
cytoskeletal mechanics and trafficking. We add a fourth arm, \emph{tau-protein
knockdown} (shRNA against MAPT), which reduces MT \emph{length}
without globally destroying the network. Under the Fr\"ohlich stress-test logic of \S 4.1.2, tau knockdown does not
test a universal length threshold. Instead, it shifts the
microtubule length distribution relative to the mode-dependent crossover
\(L_*(f_F)=c_s/(2f_F)\), reduces the number of coherently participating dimers,
and can increase the effective collective linewidth. A Fr\"ohlich-compatible
outcome would require preservation of near-\(\beta\simeq0\) linewidth scaling
together with \(N>N_{\rm crit}\) and \(\Gamma_{\rm col}\sim10^6\,{\rm s}^{-1}\).
Failure of these conditions would disfavour Fr\"ohlich as a functional channel
without excluding ordinary cytoskeletal contributions to synchronisation.
\textbf{Statistical design.} $2\times 2$ factorial (tau-KD vs scramble,
actin-KD vs scramble) with $n=10$ cultures per cell; primary endpoint
gamma PLV; secondary endpoint $\sigma$. Pre-registration of interaction
contrast recommended.
\textbf{IQI estimate:} $\text{IQI}_{\rm spec} \approx 0.90$ (electrophysiology standard, moderate $S_{\rm prep}$ variability across cultures).

\subsection{Experiment 4: Cavity spectroscopy}
\label{sec:exp4}
\textbf{Objective:} Test QED cavity model predictions \cite{Mavromatos2025}.
\textbf{Predictions.} The QED cavity model predicts a collective Rabi
splitting $\Omega_R \approx 6\THz \approx 25\meV$ (Eq.~\ref{eq:rabi_splitting}).
The corresponding wavelength splitting at 280~nm spans a range depending
on the transition dipole used: $\Delta\lambda \approx 1.6$~nm
(Eq.~\ref{eq:delta_lambda_nom}, optical-transition estimate) or
$\Delta\lambda \approx 0.89$~nm (Eq.~\ref{eq:delta_lambda_num},
conformational-dipole numerical estimate). These define the full target
window; both are below standard UV linewidths and require the strategies
detailed below.
\textbf{Challenge: linewidth and turbidity.} Standard protein UV
absorption FWHM is $\sim 20\text{--}30$~nm, far exceeding either predicted
$\Delta\lambda$. Direct resolution of a doublet in standard absorption is
therefore \textit{unlikely}. Moreover, MT solutions are highly turbid due
to their size, leading to significant scattering that can distort the
baseline and obscure weak features. We propose three detection strategies
that mitigate these issues:
\textit{(A) Second-derivative UV spectroscopy.} The second derivative
$d^2A/d\lambda^2$ of the absorption spectrum resolves sub-linewidth
structure. A doublet with separation $\Delta\lambda$ in the range
0.89--1.6~nm produces two negative lobes separated by $\Delta\lambda$ in
the second-derivative spectrum, detectable even when the parent spectrum
appears as a single peak. Scattering contributes a slowly varying
background that is largely suppressed by differentiation. Resolution:
$< 0.5$~nm. Compare polymerized MTs vs free tubulin vs denatured protein.
\textit{(B) THz/far-IR spectroscopy.} The Rabi splitting directly probes
the conformational mode at $\sim 6\THz$. Time-domain THz spectroscopy of
stabilized MT films should reveal absorption peaks absent in free
tubulin. This approach is less affected by scattering because the
wavelength is large compared to MT dimensions, and the sample can be
prepared as a thin film or in a low-scattering medium.
\textit{(C) 2D electronic spectroscopy near 280~nm.} Cross-peaks between
the Rabi-split states would appear as off-diagonal features, as in
photosynthetic complex studies \cite{Thyrhaug2018,Thyrhaug2018b}, with methodological details on polarization control and vibrational assignment provided in \cite{Thyrhaug2018b}. This technique
inherently separates homogeneous and inhomogeneous broadening and can
resolve splittings smaller than the linear absorption linewidth.
\textbf{Protocol:} (1) stabilize MTs with taxol (10~$\mu$M); (2) immobilize
on functionalized glass surface (for UV) or prepare thin films (for THz);
(3) measure with approach A, B, or C; (4) compare polymer/monomer/denatured
controls; (5) fit to doublet vs single-peak model.
\textbf{Interpretation.} Doublet or THz absorption observed: strong
evidence for cavity model, $\tau \sim \mu$s plausible. Single peak only:
cavity model falsified at predicted coupling strength, equilibrium
$\tau \sim$~ps confirmed.

% ============================================================
% FIGURE 5: SPECTRAL SIGNATURES
% ============================================================
\begin{figure*}[t]
\centering
\includegraphics[width=0.95\textwidth]{figures_final/fig2_signatures.pdf}
\caption{\textbf{Spectroscopic signatures and detection-power analysis for
ordered-water cavity hypotheses.}
Candidate ordered-water QED cavity models predict polaritonic splitting and
spectral doublet structure in UV/THz observables. Upper left: avoided-crossing
structure of the upper and lower polaritonic branches as a function of detuning.
Upper right: second-derivative UV absorption signature associated with a finite
splitting. Bottom: operational detection-power calculation using
\(\Delta{\rm BIC}\) as the decision statistic, with decisive evidence defined
by \(\Delta{\rm BIC}>10\). This figure implements the detectability layer of
the framework: a mechanism is not treated as supported merely because it is
physically possible, but only if the predicted spectral feature exceeds the
instrumental SNR and model-selection threshold.}
\label{fig:spectral_signatures}
\end{figure*}

\textbf{Why not observed previously?} Most protein UV studies use soluble
tubulin (no cavity); MT solutions have high turbidity; predicted splitting
not sought before Mavromatos 2025; standard spectral resolution
($\sim 1$~nm) marginal for the conformational-dipole estimate and
insufficient for the optical estimate.
\textbf{Feasibility:} MEDIUM. Approaches A and B available in most
biophysics labs. Cost $\sim$ \$50K--100K. Timeline 12--18 months.
\textbf{IQI estimate:} $\text{IQI}_{\rm spec} \approx 0.85$ (requires careful scattering subtraction; $S_{\rm prep}$ depends on turbidity control).

\subsection{Bayesian model selection and detection-power surface}
\label{sec:bic_roc}
To quantify the statistical evidence for a doublet vs singlet spectrum,
we employ the Bayesian Information Criterion (BIC). For a given
signal-to-noise ratio (SNR), the difference
$\Delta\text{BIC} = \text{BIC}_{\text{singlet}} -
\text{BIC}_{\text{doublet}}$ indicates the strength of evidence. Using
our current idealized synthetic fit (20 effective resolution elements,
direct doublet-vs-singlet regression), we find that decisive evidence starts
around SNR $\sim 2.6\times 10^2$ and is unambiguous by SNR $\gtrsim 10^3$,
$\Delta\text{BIC} > 10$, corresponding to ``decisive'' evidence for the
doublet model in the Jeffreys scale. This provides a quantitative
threshold for experimental design: achieving SNR $\geq 10^3$ at
the predicted splitting should be a primary instrument-specification
target.
\textbf{Full detection-power surface.} A single SNR threshold is less informative than
the full detection-probability surface. We computed, via Monte-Carlo
($10^3$ realisations per grid point), the probability
$P(\text{detect doublet}\mid\Delta\lambda,\text{SNR})$ for
$\Delta\lambda\in[0.5,2.0]$~nm and $\text{SNR}\in[10,10^4]$ (logarithmic
grid), declaring detection when $\Delta\text{BIC}>10$. Qualitative summary
of the resulting surface:
\begin{itemize}
\item Along the nominal target $\Delta\lambda=1.6$~nm: $P=0.5$ at
SNR$\approx 6\times 10^2$; $P=0.95$ at SNR$\approx 2\times 10^3$.
\item Along the numerical target $\Delta\lambda=0.89$~nm: $P=0.5$ at
SNR$\approx 2.5\times 10^3$; $P=0.95$ at SNR$\approx 8\times 10^3$.
\item Below $\Delta\lambda\simeq 0.5$~nm detection requires
SNR$>10^4$, beyond current feasibility for scattering-limited UV; the
THz approach (B) becomes mandatory.
\end{itemize}
Reframing Experiment~4 in these Neyman--Pearson detection-power terms shifts the primary instrument
specification from ``SNR $>100$ at one wavelength'' (Jeffreys threshold
on the singlet baseline) to ``SNR $\gtrsim 10^3$ across the
277--283~nm window for high-confidence detection at the
conformational-dipole target.'' This is achievable with modern
derivative-UV spectrophotometers but requires careful scattering
subtraction.

\subsection{Experiment 5: Frequency- and length-resolved THz linewidth
as a Fr\"ohlich stress test}
\label{sec:exp5}
The corrected Fr\"ohlich analysis shifts the role of this experiment. It is
not designed to confirm a generic length-gating rescue mechanism. For
sub-THz modes, the crossover length
\[
L_*(f_F)=\frac{\pi c_s}{\omega_F}
=\frac{c_s}{2 f_F}
\]
is only $\sim10$--$100$ nm for $f_F=10^{10}$--$10^{11}$ Hz, so ordinary
polymerized MTs already satisfy $L\gg L_*$. The decisive quantities are
therefore the absolute collective damping $\Gamma_{\rm col}$, the scaling
exponent $\beta$, and the number of coherently participating dimers.

\textbf{Protocol.} Polymerize MTs under controlled nucleation to obtain
length distributions peaked at
$L\in\{0.5,1,2,5,10,20,50\}\,\mu$m where feasible, validated by cryo-EM
or high-resolution fluorescence sizing. Measure time-domain THz/far-IR
absorption over the $0.01$--$1$ THz window. Primary endpoints are:
(i) the absolute linewidth $\Gamma_{\rm col}(f_F,L)$;
(ii) the scaling exponent
\[
\beta(f_F)=\frac{d\log\Delta\omega_F}{d\log N};
\]
and (iii) whether the measured pair $(N,\Gamma_{\rm col})$ satisfies the
Fr\"ohlich threshold condition
\[
N>
6.4\times10^4
\left(\frac{\Gamma_{\rm col}}{10^6\,{\rm s}^{-1}}\right)
\left(\frac{r_{\rm pump}}{1\,{\rm s}^{-1}}\right)^{-1}.
\]

\textbf{Decision logic.}
\begin{itemize}
\item If $\Gamma_{\rm col}\gg10^6\,{\rm s}^{-1}$, Fr\"ohlich condensation is
below threshold even if $\beta\simeq0$.
\item If $\beta\gtrsim0.5$, the required coherent participation rapidly
becomes biologically implausible.
\item If $\beta\simeq0$ and $\Gamma_{\rm col}\lesssim10^6\,{\rm s}^{-1}$,
Fr\"ohlich pumping survives as a ps-scale non-equilibrium spectroscopic
mechanism, but still does not close the neural functional gap without a
large independently measured amplification cascade.
\item If no narrow collective mode is observed in the predicted band,
Fr\"ohlich condensation is strongly disfavoured in biologically relevant MT
preparations.
\end{itemize}

\textbf{Interpretation.} A positive result in Experiment~5 would not establish
neural quantum functionality. It would only rescue Fr\"ohlich pumping as a
candidate collective vibrational phenomenon. A negative result would sharply
constrain or effectively falsify Fr\"ohlich condensation in microtubules under
the tested biological conditions.

\subsection{Experiment 6: Length-resolved linewidth scaling for subradiance
characterisation}
\label{sec:exp6}
\textbf{Objective:} Constrain the $(\tau_A, p)$ region of
Proposition~\ref{prop:subradiance_threshold} and reconcile the Babcock
et al.\ 10-s subradiant lifetime with the symmetry-breaking mixing rate
$\Gamma_{\text{mix}} \sim 10^{10}\text{--}10^{11}$~Hz.
\textbf{Physical target.} The critical threshold
$\tau^*(p) = 2.5\times 10^{-9}/p$~s (Eq.~\ref{eq:subradiance_threshold})
can be traversed experimentally by varying MT length, since the geometric
symmetry underlying subradiance degrades with increased translational
disorder at shorter lengths. Specifically:
\begin{itemize}
\item In long, well-ordered MTs ($L \gtrsim 5\,\mu$m) the tryptophan
network has higher translational periodicity, lower effective disorder
$\delta\omega/\omega$, and reduced $\Gamma_{\text{mix}}$.
\item In short or nocodazole-treated MTs the disorder is high and
$\Gamma_{\text{mix}}$ dominates, driving the system toward the
disorder-dominated row of the table in \S\ref{sec:subradiance}.
\end{itemize}
The predicted signature is a \emph{length-dependent narrowing} of the
subradiant UV emission linewidth $\Delta\nu_{\text{sub}}$:
\begin{equation}
\Delta\nu_{\text{sub}}(L) \propto
\Gamma_{\text{mix}}(L) \approx
\left(\frac{\delta\omega(L)}{\omega}\right)^2 \Gamma_\phi,
\end{equation}
where $\delta\omega(L)/\omega$ decreases with $L$ as the tryptophan lattice
becomes more ordered. A plot of $\Delta\nu_{\text{sub}}$ vs $\log L$
provides a direct upper bound on $\tau_A^{\text{effective}}$ and,
combined with fluorescence-lifetime resolved spectra, maps the
experimentally realised $(p, \tau_A)$ point onto
Proposition~\ref{prop:subradiance_threshold}.
\textbf{Protocol.}
(1)~Polymerise MTs to length distributions at
$L \in \{1, 2, 5, 10, 20\}\,\mu$m (cryo-EM verified).
(2)~UV fluorescence spectroscopy (280~nm excitation) with picosecond
time-correlated single-photon counting to resolve the fast superradiant
($\tau_S \sim 100\fs$--ps) and slow subradiant components.
(3)~Extract the subradiant component lifetime $\tau_A(L)$ and emission
linewidth $\Delta\nu_{\text{sub}}(L)$ per length bin.
(4)~Compare taxol-stabilised vs nocodazole-treated vs free-tubulin controls.
(5)~Test disorder-reduction hypothesis: high-ionic-strength buffer
(reduced Debye shielding, more ordered lateral contacts)
vs low-ionic-strength.
\textbf{Primary endpoint:} the slope $d\log\tau_A / d\log L$.
\begin{itemize}
\item Steep positive slope ($\tau_A \propto L^\alpha$, $\alpha > 1$):
subradiance is geometry-gated; long MTs are in the
$p > 2.5\times 10^{-3}$ region ($\mathcal{U} \ge 1$ marginal).
\item Flat slope with $\tau_A$ near 10~s for all lengths:
subradiant protection is intrinsic rather than length-gated;
Babcock et al.\ is reproduced, supports
$p \gg 2.5\times 10^{-9}$ and $\mathcal{U} \gg 1$.
\item Short-length $\tau_A \sim 10\ps$ (disorder-dominated):
no $\mathcal{U} \ge 1$ crossing, subradiance falsified as
a functional mechanism.
\end{itemize}
\textbf{Reconciliation with mixing rate.} If $\tau_A \sim 10$~s is
reproduced at any length, then $\Gamma_{\text{mix}} \sim 10^{10}$~Hz must
be reconciled with the long lifetime. Possible resolutions (to be tested):
(a)~the relevant disorder parameter is $\delta\omega \ll \omega$ due to
topological protection in helical MT geometry; (b)~dynamical symmetrisation
by the QED cavity field (links to Experiment~4); (c)~coupling to dark
vibrationally excited states that buffer the symmetry-breaking perturbation.
This reconciliation constitutes an explicit theoretical target for the
v2 structured-bath HEOM extension.
\textbf{Statistical design.} $n = 5$ MT preparations per length class,
$n = 3$ independent fluorescence measurements per preparation.
Power: with the expected factor $>10^{10}$ between subradiant and
superradiant lifetime, even modest $n$ provides $>90\%$ power to
detect the subradiant component above noise.
Cost $\sim$\$80\,K; timeline 12--18 months.
Feasibility: MEDIUM-HIGH (requires UV ultrafast spectroscopy; comparable
to Babcock et al.\ setup).
\textbf{IQI estimate:} $\text{IQI}_{\rm spec} \approx 0.94$ (high-precision TCSPC, strong $S_{\rm prep}$ via cryo-EM validation).

% ============================================================
% DISCUSSION
% ============================================================
\section{Discussion}
\subsection{Principal findings}
(1) The calibrated Ohmic equilibrium pure-dephasing baseline gives
\(T_2^*\approx39\fs\) at 310~K, while structured-bath excitonic benchmarks
remain in the ps regime. Both scales are many orders of magnitude below
neural synchronization times.
(2) The quantum coherence utility $\mathcal{U}_{\rm phys} < 1$ for all equilibrium
cooperative mechanisms (Definition~\ref{def:utility}); quantum discord
persists to $\sim 2\tau_c$ without changing this conclusion
(\S\ref{sec:discord}).
(3) The Holstein mapping places tubulin in the \emph{small-polaron regime}
($J/\hbar\Omega_0 \sim 10^{-3}$), suppressing coherent inter-dimer
electronic transport by $(J/\hbar\Omega_0)^2 \sim 10^{-6}$ relative to
FMO. Critically, this is independent of the 2D-IR cross-peak amplitude,
which is set by $S \approx 1$ and predicts a cross/diagonal ratio of
$\sim 3$--$5\%$---detectable but not diagnostic of functional ENAQT.
Fr\"ohlich pumping represents the specific pathway to enter the ENAQT
window, testable via the length-dependent cross-peak increase in
Experiments~1 and~5.
(4) Fr\"ohlich condensation is no longer treated as a presumed bridge to
functionality. For sub-THz modes, the form-factor crossover
\(L_*(f_F)=c_s/(2f_F)\) is typically in the \(10\)--\(100\) nm range, so
micrometre-scale microtubules can in principle enter the favourable
\(\beta\simeq0\) geometric regime. The bottleneck instead becomes the
simultaneous requirement of large coherent participation
(\(N_{\rm crit}\sim6.4\times10^4\) for
\(\Gamma_{\rm col}=10^6\,{\rm s}^{-1}\)) and exceptionally weak collective
damping. Fr\"ohlich pumping is therefore a stress-tested but currently
disfavoured functional channel, whereas QED cavity protection and geometric
subradiance remain the only candidate mechanisms that can approach
\(\mathcal{U}\gtrsim1\) without an extreme downstream gain.
(5) The quantum-to-classical transduction bandwidth $B_{\text{QC}}
\lesssim 10^4$~Hz (ion channel kinetics) limits the effective information
gain to $\Delta R_{\text{eff}} \lesssim 10^4$~bits/s under any equilibrium
scenario; $B_{\text{QC}}$ unifies the utility framework (\S 3.4) and the
information-theoretic analysis (\S 5.5) and is the physical quantity
targeted by Experiment~3 (Eq.~\ref{eq:U_BQC}).
(6) All current neural synchronization phenomena are explained by
classical mechanisms (\S 5.4--5.5).
(7) UV superradiance \cite{Babcock2024}, anesthetic-MT interactions
\cite{Khan2024,Kalra2023}, and MT resonances \cite{Saxena2020,Singh2021}
support further experimental investigation.

\subsection{Reconciliation of conflicting predictions}
Equilibrium models \cite{Tegmark2000,Hagan2002,Xiang2019} yield
$\tau \sim 10\fs\text{--}10\ps$. Non-equilibrium models diverge: Fr\"ohlich
(optimistic) $\tau \sim 10\text{--}100\ps$; Fr\"ohlich (Reimers critique)
inaccessible \cite{Reimers2009}; cavity QED $\tau \sim 0.1\text{--}1~\mu$s
\cite{Mavromatos2025}; subradiance up to 10~s \cite{Babcock2024}.
Key distinctions: (a) equilibrium vs driven systems---Corollary~\ref{cor:ceiling}
assumes thermal equilibrium, violated in living systems dissipating
$\sim 100$~W metabolic power; (b) geometry---Reimers \cite{Reimers2009}
analyzed 1D chains, MTs are 2D cylindrical lattices, and the $\beta$
analysis (\S 4.1.2) shows that geometry+length dictate which regime
applies; (c) nonlinearity---Reimers used linearized dynamics, Fr\"ohlich
condensation requires nonlinear dipole interactions \cite{Pokorn2011};
(d) collective vs independent modes---\S 3.2 shows independent baths give
$\tau_N = \tau_1/N$, while Fr\"ohlich condensation involves a single
collective mode avoiding this scaling.

\subsection{Comparison with FMO: lessons from updated understanding}
The FMO case initially appeared to support long electronic coherence.
Careful analysis revealed electronic coherence is ultrafast ($\sim 60\fs$)
\cite{Duan2017,Thyrhaug2018}, vibrational coherences persist longer
($\sim$~ps) \cite{Thyrhaug2018}, and functional energy transfer uses
vibronic mechanisms plus correlated environmental fluctuations
\cite{Dean2017,Cao2020,Chin2013}. The Holstein mapping of
\S\ref{sec:holstein} therefore remains a useful theoretical organiser, but
the calibrated tubulin release described here does not by itself place the
system in the same numerical coherence window as FMO vibronics; any such
claim requires structured-bath or driven extensions beyond the present code.

\subsection{Limitations}
The two-level approximation neglects tubulin's complex conformational
landscape; qudit models \cite{Mavromatos2025} are more realistic. Spectral
density parameters are generic protein values (\S 2.1.1); tubulin-specific
MD simulations should refine $\eta$ and $\omega_c$. The Fr\"ohlich
$\beta$-estimate (\S 4.1.2) is a linear-response result; nonlinear
self-trapping could reduce $\beta$ further, slightly extending the
optimistic regime. The QED cavity model relies on poorly characterized
ordered water structure; the two estimates of $\Delta\lambda$
(Eqs.~\ref{eq:delta_lambda_nom}--\ref{eq:delta_lambda_num}) reflect
genuine uncertainty in the relevant transition dipole. The subradiant
scenario (\S\ref{sec:subradiance}) requires a full quantitative
reconciliation of the 10-s Babcock lifetime with the
$\sim$10-ps symmetry-breaking mixing rate; we flag this as priority
theoretical work.

\subsection{Falsifiability of competing non-equilibrium models}
The disagreement between Reimers et al.\ and Mavromatos et al.\ is often
presented as an irresolvable theoretical dispute. Our experimental
proposals (\S\ref{sec:experimental_discrimination}) turn it into an empirical question.
Specifically:
\begin{itemize}
\item If Experiment 4 detects the predicted Rabi splitting in the
range $\Delta\lambda \approx 0.89\text{--}1.6$~nm, or a THz absorption at
$\sim 6$~THz, the QED cavity model is strongly supported. Conversely, the
absence of any splitting at SNR$>10^4$ falsifies the cavity model at the
predicted coupling strengths, lending weight to Reimers' conclusion.
\item If Experiment 5 finds $\beta\lesssim 0.1$ with length-dependent
THz linewidth, Fr\"ohlich condensation is supported; $\beta\gtrsim 0.9$
falsifies it; intermediate values narrow the viable parameter region.
\item If Experiment 3 shows specific tau-KD suppression of gamma
synchrony beyond classical structural effects, a length-dependent
quantum contribution is implicated, and $B_{\text{QC}}$ is bounded from
above (\S 5.5, Eq.~\ref{eq:U_BQC}).
\item If Experiment 6 finds $\tau_A(L) \propto L^\alpha$ with $\alpha>1$,
geometric subradiance is length-gated; if $\tau_A \sim 10$~s is reproduced
at all lengths, the Babcock et al.\ observation is generic and
$\mathcal{U}_{\text{sub}} \gg 1$ (Proposition~\ref{prop:subradiance_threshold}).
\end{itemize}
Thus, the proposed experiments provide clear, falsifiable tests that
can resolve the long-standing controversy.

\subsection{Relation to Orch OR}
Orch OR \cite{Penrose1994,Hameroff2014} provides historical motivation for
considering microtubules as possible quantum substrates, but it is not assumed
or tested in the present work. Our analysis addresses a narrower and more
operational biological-physics question: whether tubulin can host molecular
quantum dynamics with experimentally distinguishable signatures on functionally
relevant timescales. Under equilibrium assumptions the answer is negative
($\tau\sim$~fs--ps, $\mathcal{U}\ll1$). Non-equilibrium mechanisms can become
formally compatible only if they supply strong amplification
($\tau\sim\mu$s--s for the most favourable subradiant estimates,
$\mathcal{U}\sim1$--$40$) and, in several cases, cascading
quantum$\to$classical amplification ($\mathcal{K}\sim10^3$--$10^4$). Kerskens
and L\'opez P\'erez \cite{Kerskens2022} reported MRI evidence interpreted as
non-classical brain function, whereas Warren \cite{Warren2023} provided a
critical rebuttal. The present manuscript remains independent of that dispute:
it supplies constraints and experimental discriminants rather than an Orch OR
endorsement.

\subsection{Computational framework}
\label{sec:code}
All numerical results presented in this release are generated by the
open-source project \texttt{KwanTube}, available at
\url{https://github.com/FacundoFirmenich/KwanTube}. The repository contains
the manuscript-oriented entry points, validation ledgers, figure-generation
scripts, and Bayesian contraction summaries used for this bioRxiv v1 release. The code
implements the calibrated rate-model baseline described in \S\S 2--4,
performs multi-temperature inversion to constrain $\eta$ and $\omega_c$
(yielding $T_2^* = 39.1\fs$ at 310~K; \S 2.3), computes Sobol-style
sensitivity indices for the current four-parameter release
($\eta,\omega_c,T,f_{\text{prot}}$), and provides a synthetic BIC model
selection tool for distinguishing cavity doublets from singlets
(\S\ref{sec:bic_roc}). A complete validation report is included
in the supplementary information and in the repository as
\texttt{validation\_report.json}. All figures can be reproduced by
running \texttt{python src/generate\_paper\_figures.py}, which writes the
manuscript-ready assets to \texttt{figures\_final/}
(\texttt{fig1\_landscape}, \texttt{fig2\_signatures},
\texttt{fig3\_frohlich}, \texttt{fig4\_scaling}; PDF and PNG).
\textbf{Software architecture, solvers and regression tests.}
\emph{(i) Open-system engine.} The public release exposes a
phenomenological secular/Markovian rate model, not a production HEOM stack;
HEOM and structured-bath benchmarks are integrated as validation modules.
\emph{(ii) Integrals.} The vibrational-dephasing integral is evaluated with
\texttt{scipy.integrate.quad}, using an analytic low-frequency completion and
an upper cutoff $50\,\omega_c$.
\emph{(iii) Inversion.} The multi-temperature parameter inversion uses
\texttt{scipy.optimize.least\_squares} (Trust Region Reflective) with
multi-start initialisation and variance-scaled residuals.
\emph{(iv) Sensitivity indices.} The current release computes Jansen-style
first-order and total-order indices for four parameters
($\eta,\omega_c,T,f_{\text{prot}}$) with $N=2000$ base samples.
\emph{(v) Regression tests.} The test suite
(\texttt{test\_physics.py}, \texttt{test\_inversion.py}) verifies: positivity
of all rates; the FDT bound $\tau\le\hbar/(\alpha\kB T)$; near-consistency
between numeric and analytic vibrational rates; ionic-mass sensitivity; and
synthetic recovery of $(\eta,\omega_c,\Delta\epsilon)$ at the calibrated
tolerance level. A CI workflow on
every commit re-runs these tests and rebuilds all figures. The repository
state at publication is archived via Zenodo (DOI: 10.5281/zenodo.19744600)
to guarantee computational reproducibility.

% ============================================================
% FIGURES (SUPPLEMENTARY REFERENCES)
% ============================================================


\subsection{Generalization beyond tubulin: a framework for biological quantum relevance}
The equilibrium-to-functionality framework developed here is intentionally
general. While tubulin provides a stringent test case due to its structural
order and biological centrality, the formalism applies to any biomolecular
system where quantum dynamics might influence function across timescale
separations. For chromatin, the gap spans electronic coherence (fs) to
epigenetic regulation (minutes); for ion channels, it spans tunnelling
events (fs--ps) to gating dynamics (ms). The utility criterion
$\mathcal{U}_{\rm phys} = \mathcal{K}\tau_{\text{coh}}/\tau_{\text{func}}$ provides
a universal dimensionless metric: $\mathcal{U}\gtrsim 1$ defines the
necessary condition for functional quantum influence, regardless of the
specific molecular substrate. The three non-equilibrium mechanisms
analyzed---metabolic pumping, environmental cavity protection, and
geometric decoherence-free subspaces---represent generic physical pathways
that could, in principle, bridge the gap in diverse biological contexts.
Future work should apply this framework systematically to other candidate
systems (e.g., voltage-gated ion channels, chromatin remodeling complexes,
olfactory receptors) to map the landscape of biologically plausible quantum
effects. This generalization transforms the question from ``is quantum
mechanics relevant in biology?'' to the more precise ``under what
non-equilibrium conditions, and in which specific biological architectures,
can molecular quantum dynamics achieve functional relevance?''.

\section{Conclusions}
We have presented a comprehensive analysis of quantum coherence mechanisms
in microtubules with four main contributions: a calibrated open-system
baseline for equilibrium decoherence, an
FDT-rigorous quantum coherence utility framework (\S 3.4, including
quantum-discord refinement and comparison with Giovannetti-2004 /
Chin-2013), a perturbative scaling analysis of collective damping that
renders Fr\"ohlich condensation length-gated (\S 4.1.2), and the
integration of geometric subradiance as a third non-equilibrium channel
(\S\ref{sec:subradiance}), together with a Holstein vibronic analogue that
formalises the FMO--tubulin analogy (\S\ref{sec:holstein}).
\textbf{Equilibrium:} Vibrational dephasing limits coherence to
$T_2^* \approx 39\fs$ in the calibrated Ohmic release at 310~K. The strict FDT pure-dephasing ceiling remains sub-ps, and even favourable
structured-bath benchmarks remain confined to the fs--ps regime; no
cooperative equilibrium mechanism circumvents this bound.
\textbf{Non-equilibrium:} Fr\"ohlich pumping can in principle generate
driven collective vibrational behaviour, but the corrected energetic and
geometric accounting makes it a strongly constrained mechanism. For
sub-THz modes, ordinary MTs already exceed the form-factor crossover length;
the decisive bottlenecks are instead low absolute damping, coherent
participation above $N_{\rm crit}$, and downstream amplification. Even in
the favourable case, the expected coherence time remains
$\tau\sim10$--$100$ ps, requiring $\mathcal{K}_{\rm req}\sim10^8$--$10^9$
to influence millisecond neural function. Fr\"ohlich therefore shifts from
a presumed functional bridge to a falsifiable but currently disfavoured
functional channel. QED cavity models predict $\tau\sim0.1$--$1~\mu$s and
remain testable by UV/THz cavity spectroscopy. Geometric subradiance,
reported in related tryptophan mega-networks \cite{Babcock2024}, reaches
lifetimes up to 10 s but requires reconciliation with vibrational and ionic
mixing and with functional transduction from optically dark states.
\textbf{Classical sufficiency:} ING networks, gap junctions, ephaptic
coupling, and CTC quantitatively explain all observed neural
synchronization (\S 5.4). Information-theoretic analysis confirms no rate
advantage for quantum coherence under equilibrium assumptions (\S 5.5).
\textbf{Experimental priorities:} (1) Cavity spectroscopy---second-derivative
UV or THz---for direct test of QED cavity model, with detection-power design
specifying SNR $\gtrsim 10^3$ as a conservative target and idealized BIC
crossing already in the few-$10^2$ range; (2) temperature-dependent 2D-IR with cross/diagonal amplitude
ratio of $\sim 3$--$5\%$ as the equilibrium target (vibronic coupling,
$S$-controlled) and a \textit{length-dependent increase above this baseline}
as the ENAQT indicator (transport efficiency, $J/\hbar\Omega_0$-controlled);
the two quantities are physically orthogonal and must be reported separately; (3) kinetic isotope effects for tunneling;
(4) tau-knockdown-controlled MT disruption with synchronization precision
measurements, directly constraining $B_{\text{QC}}$; (5) length-resolved
THz linewidth as direct $\beta$-falsification test for Fr\"ohlich;
(6) length-resolved UV subradiant linewidth to map the $(\tau_A, p)$
region and reconcile the Babcock 10-s lifetime with the mixing rate.
The amplification gap from picoseconds to milliseconds remains a hard
constraint, not a rhetorical one. Our utility framework sharpens that
constraint quantitatively and converts it into a testable experimental
agenda. In that sense, the contribution of this manuscript is twofold:
it narrows the plausible mechanism space under explicit assumptions, and
it specifies concrete falsification routes for cavity, Fr\"ohlich, and
subradiant scenarios without relying on overextended interpretative claims.

% ============================================================
% ETHICS, ACKNOWLEDGMENTS, AUTHOR CONTRIBUTIONS
% ============================================================
\section*{Ethics and Conflict of Interest}
The authors declare no financial or personal conflicts of interest.

\section*{Acknowledgments}
Numerical analysis and visualization were performed using the \textbf{Python} scientific stack, including \textbf{NumPy}, \textbf{SciPy}, \textbf{Matplotlib}, and \textbf{Pandas}. We thank \textbf{Marc Mart\'inez Sep\'ulveda}, \textbf{Hernando Rub\'en Ch\'avez Fern\'andez}, \textbf{Camila Paz Bravo}, and \textbf{Pablo Cingolani}, along with the entire team at \textbf{CEDESUR}. We also acknowledge \textbf{Ernesto Villanueva} and \textbf{Juan Pastor Gonz\'alez} (UNAJ) for institutional support.

\section*{Data and code availability}

The canonical computational entry point for reproducing the manuscript-level
results is the Python workflow
\texttt{scripts/reproduce\_paper\_results.py}. The platform-specific launchers
provided in the repository are convenience wrappers around this canonical
Python entry point; they do not replace it.

\begin{itemize}
\item \textbf{Primary repository:}
\url{https://github.com/FacundoFirmenich/KwanTube}
(GPL-3.0 License).
\item \textbf{Permanent archive:}
\href{https://doi.org/10.5281/zenodo.19744600}{10.5281/zenodo.19744600}.
\item \textbf{Canonical command:}
\texttt{python scripts/reproduce\_paper\_results.py --mode paper}.
\item \textbf{Convenience launchers:}
\texttt{reproduce\_paper.bat} for Windows,
\texttt{reproduce\_paper.sh} for Linux/macOS terminals,
\texttt{reproduce\_paper.command} for macOS Finder, and
\texttt{reproduce\_paper\_linux.desktop} for Linux desktop environments that
allow trusted executable launchers.
\end{itemize}

The launchers create or activate a local Python environment, install the
declared dependencies, run the canonical reproducibility workflow, and rebuild
the manuscript figures and validation outputs. Where a LaTeX installation is
available, the workflow can also be used to rebuild the manuscript PDF. The
Python workflow remains the authoritative reproducibility layer.

\section*{Hardware disclosure}

The manuscript-level Python workflow was designed and tested on commodity
hardware, including a mid-range 2016 laptop with an Intel i5-6200U CPU,
16 GB DDR3 RAM, SSD storage, and 128 MB dedicated VRAM. Baseline figure
generation, rate-model calculations, validation-table regeneration, and
synthetic detection-power analyses do not require high-end HPC access.
Optional long-window HEOM production extensions are treated separately as
staged computational extensions and may require substantially longer runtimes
depending on truncation depth, bath expansion, and system size.

\section*{Author contributions}

\textbf{Facundo Firmenich} led conceptualization, theoretical development,
methodology, formal analysis, software implementation, validation,
visualization, project administration, and original-draft writing.
\textbf{Pau Firmenich} contributed to reproducibility testing, computational
workflow review, validation checks, and editorial quality control.
\textbf{Le\'on Firmenich} contributed to independent review of figures,
tables, explanatory consistency, and editorial quality control. All authors
reviewed and approved the final manuscript and associated software release.

\section*{AI-assisted development disclosure}
\label{sec:ai_disclosure}

Large language models, including ChatGPT, Qwen, DeepSeek, GLM, and Gemini,
were used as assistive tools for drafting support, code prototyping,
debugging, translation, refactoring, and consistency checking. Wolfram Alpha
was used for symbolic and numerical spot checks. All scientific claims,
equations, numerical outputs, code behaviour, interpretations, and final
editorial decisions were reviewed, validated, and approved by the human
authors. The AI systems were used as research-assistance tools, not as
autonomous scientific agents.


% ============================================================
% APPENDICES AND SUPPLEMENTARY FIGURES
% ============================================================
\appendix
\section{Supplementary material}
Detailed derivations of $\mathcal{F}_{\rm info}(\text{SNR})$, full Fisher barrier calibration, spectral posterior predictive checks (PPC), dual KL+$P_{\rm boundary}$ diagnostics, and complete L1--L4 validation ledgers are provided in the supplementary information to maintain main-text flow while preserving full methodological transparency.

% Supplementary figures referenced in main text
\begin{figure*}[t]
\centering
\includegraphics[width=0.95\textwidth]{figures_final/redfield_1JFF.pdf}
\caption{\textbf{Redfield baseline dynamics for the full 1JFF system.}
Perturbative Redfield trajectory for the full 1JFF tubulin-inspired system.
Left: site-resolved population redistribution over the short-time window.
Right: \(L_1\)-norm coherence. This figure provides the perturbative baseline
against which the HEOM short-window benchmark is compared. It is used as a
model-comparison diagnostic and not as a numerically exact reference.}
\label{fig:supp_redfield_1jff}
\end{figure*}

\begin{figure*}[t]
\centering
\includegraphics[width=0.95\textwidth]{figures_final/redfield_6DPU.pdf}
\caption{\textbf{Redfield baseline dynamics for the 6DPU fragment.}
Site-resolved Redfield population dynamics and \(L_1\)-norm coherence for the
6DPU fragment. The many-site legend is retained to preserve traceability of the
fragment-level diagnostic, but the figure is treated as a supplementary
validation artefact rather than a main-text result. Together with the 1JFF
baseline, it documents the perturbative reference layer used in the
HEOM/Redfield validation ledger.}
\label{fig:supp_redfield_6dpu}
\end{figure*}

\begin{figure*}[t]
\centering
\includegraphics[width=0.95\textwidth]{figures_final/prior_sensitivity.pdf}
\caption{\textbf{Prior-sensitivity analysis for evidence ratios.}
Sensitivity of \(BF_{10}\) to prior width for the Babcock and Kalra evidence
layers. Horizontal dashed lines mark conventional strong and decisive evidence
thresholds. The Babcock analysis remains in the strong-to-decisive regime over
a broad prior-width interval, while the Kalra layer is more prior-sensitive and
approaches the strong-evidence threshold for broad priors. The figure supports
the use of study-specific Bayesian evidence as a diagnostic layer while
discouraging overinterpretation of any single pooled Bayes factor.}
\label{fig:supp_prior_sensitivity}
\end{figure*}

\begin{figure*}[t]
\centering
\includegraphics[width=0.85\textwidth]{figures_final/sbc_calibration_ns.pdf}
\caption{\textbf{Simulation-based calibration of the nested-sampling engine.}
Rank histogram for simulation-based calibration of the nested-sampling evidence
engine. The dashed line denotes the expected uniform rank distribution under
calibrated posterior recovery. The reported \(p=0.8009\) is consistent with no
detectable deviation from rank uniformity in this calibration test, supporting
use of the engine for evidence diagnostics within the validation ledger.}
\label{fig:supp_sbc_ns}
\end{figure*}

% ============================================================
% BIBLIOGRAPHY
% ============================================================
\begin{thebibliography}{99}
\bibitem{Engel2007} G.~S.~Engel \textit{et al.}, \textit{Nature}
\textbf{446}, 782 (2007).
\bibitem{Panitchayangkoon2010} G.~Panitchayangkoon \textit{et al.},
\textit{Proc.\ Natl.\ Acad.\ Sci.\ USA} \textbf{107}, 12766 (2010).
\bibitem{Duan2017} H.-G.~Duan \textit{et al.},
\textit{Proc.\ Natl.\ Acad.\ Sci.\ USA} \textbf{114}, 8493 (2017).
\bibitem{Thyrhaug2018} E.~Thyrhaug \textit{et al.},
\textit{Nat.\ Chem.} \textbf{10}, 780 (2018).
\bibitem{Thyrhaug2018b} E.~Thyrhaug \textit{et al.},
\textit{J.\ Phys.\ Chem.\ Lett.} \textbf{9}, 6412 (2018).
\bibitem{Tiwari2013} V.~Tiwari, W.~K.~Peters, and D.~M.~Jonas,
\textit{Proc.\ Natl.\ Acad.\ Sci.\ USA} \textbf{110}, 1203 (2013).
\bibitem{Christensson2012} N.~Christensson \textit{et al.},
\textit{J.\ Phys.\ Chem.\ B} \textbf{116}, 7449 (2012).
\bibitem{Tempelaar2014} R.~Tempelaar, T.~L.~C.~Jansen, and J.~Knoester,
\textit{J.\ Phys.\ Chem.\ B} \textbf{118}, 12865 (2014).
\bibitem{Dean2017} J.~C.~Dean and G.~D.~Scholes,
\textit{Acc.\ Chem.\ Res.} \textbf{50}, 2746 (2017).
\bibitem{Plenio2008} M.~B.~Plenio and S.~F.~Huelga,
\textit{New J.\ Phys.} \textbf{10}, 113019 (2008).
\bibitem{Cao2020} J.~Cao \textit{et al.},
\textit{Sci.\ Adv.} \textbf{6}, eaaz4888 (2020).
\bibitem{Chin2013} A.~W.~Chin, J.~Prior, R.~Rosenbach, F.~Caycedo-Soler,
S.~F.~Huelga, and M.~B.~Plenio,
\textit{Nat.\ Phys.} \textbf{9}, 113 (2013).
\bibitem{Ritz2000} T.~Ritz, S.~Adem, and K.~Schulten,
\textit{Biophys.\ J.} \textbf{78}, 707 (2000).
\bibitem{Maeda2012} K.~Maeda \textit{et al.},
\textit{Proc.\ Natl.\ Acad.\ Sci.\ USA} \textbf{109}, 4774 (2012).
\bibitem{Gauger2011} E.~M.~Gauger \textit{et al.},
\textit{Phys.\ Rev.\ Lett.} \textbf{106}, 040503 (2011).
\bibitem{Scrutton1999} N.~S.~Scrutton, J.~Basran, and M.~J.~Sutcliffe,
\textit{Eur.\ J.\ Biochem.} \textbf{264}, 666 (1999).
\bibitem{Klinman2013} J.~P.~Klinman and A.~Kohen,
\textit{Annu.\ Rev.\ Biochem.} \textbf{82}, 471 (2013).
\bibitem{Babcock2024} N.~S.~Babcock \textit{et al.},
\textit{J.\ Phys.\ Chem.\ B} \textbf{128}, 4035 (2024).
\bibitem{Celardo2019} G.~L.~Celardo, A.~Ferretti, and R.~Kaiser,
\textit{New J.\ Phys.} \textbf{21}, 023012 (2019).
\bibitem{Dicke1954} R.~H.~Dicke,
\textit{Phys.\ Rev.} \textbf{93}, 99 (1954).
\bibitem{Gross1982} M.~Gross and S.~Haroche,
\textit{Phys.\ Rep.} \textbf{93}, 301 (1982).
\bibitem{Mavromatos2025} N.~E.~Mavromatos, A.~Mershin, and D.~V.~Nanopoulos,
\textit{Eur.\ Phys.\ J.\ Plus} \textbf{140}, 1116 (2025).
\bibitem{Mavromatos1997} N.~E.~Mavromatos and D.~V.~Nanopoulos,
\textit{Int.\ J.\ Mod.\ Phys.\ B} \textbf{11}, 851 (1997).
\bibitem{Mavromatos2002} N.~E.~Mavromatos, A.~Mershin, and D.~V.~Nanopoulos,
\textit{Int.\ J.\ Mod.\ Phys.\ B} \textbf{16}, 3623 (2002).
\bibitem{DelGiudice1988} E.~Del~Giudice, G.~Preparata, and G.~Vitiello,
\textit{Phys.\ Rev.\ Lett.} \textbf{61}, 1085 (1988).
\bibitem{Sahu2013} S.~Sahu \textit{et al.},
\textit{Biosens.\ Bioelectron.} \textbf{47}, 141 (2013).
\bibitem{Tuszynski1995} J.~A.~Tuszy\'nski \textit{et al.},
\textit{J.\ Theor.\ Biol.} \textbf{174}, 371 (1995).
\bibitem{Frohlich1968} H.~Fr\"ohlich,
\textit{Phys.\ Lett.\ A} \textbf{26}, 402 (1968).
\bibitem{Frohlich1977} H.~Fr\"ohlich,
\textit{Riv.\ Nuovo Cimento} \textbf{7}, 399 (1977).
\bibitem{Reimers2009} J.~R.~Reimers \textit{et al.},
\textit{Proc.\ Natl.\ Acad.\ Sci.\ USA} \textbf{106}, 4219 (2009).
\bibitem{Pokorn2011} J.~Pokorn\'y, J.~Ha\v{s}ek, and F.~Jel\'inek,
\textit{J.\ Phys.:\ Conf.\ Ser.} \textbf{329}, 012013 (2011).
\bibitem{Salari2011} V.~Salari \textit{et al.},
\textit{J.\ Phys.:\ Conf.\ Ser.} \textbf{306}, 012075 (2011).
\bibitem{Khan2024} S.~Khan \textit{et al.},
\textit{eNeuro} \textbf{11}, ENEURO.0532-23.2024 (2024).
\bibitem{Kalra2023} A.~P.~Kalra \textit{et al.},
\textit{ACS Cent.\ Sci.} \textbf{9}, 352 (2023).
\bibitem{Wiest2025} M.~C.~Wiest,
\textit{Neurosci.\ Conscious.} \textbf{2025}, niaf011 (2025).
\bibitem{Saxena2020} K.~Saxena \textit{et al.},
\textit{Fractal Fract.} \textbf{4}, 11 (2020).
\bibitem{Singh2021} P.~Singh \textit{et al.},
\textit{J.\ Neurophysiol.} \textbf{125}, 2107 (2021).
\bibitem{Tegmark2000} M.~Tegmark,
\textit{Phys.\ Rev.\ E} \textbf{61}, 4194 (2000).
\bibitem{Hagan2002} S.~Hagan, S.~R.~Hameroff, and J.~A.~Tuszy\'nski,
\textit{Phys.\ Rev.\ E} \textbf{65}, 061901 (2002).
\bibitem{May2011} V.~May and O.~K\"uhn, \textit{Charge and Energy Transfer
Dynamics in Molecular Systems}, 3rd ed.\ (Wiley-VCH, 2011).
\bibitem{Rego2003} L.~G.~C.~Rego \textit{et al.},
\textit{J.\ Am.\ Chem.\ Soc.} \textbf{125}, 7989 (2003).
\bibitem{Xiang2019} Z.~Xiang, C.~Tang, and L.~Xin,
\textit{Chin.\ Phys.\ B} \textbf{28}, 048701 (2019).
\bibitem{Eaves2005} J.~D.~Eaves \textit{et al.},
\textit{Proc.\ Natl.\ Acad.\ Sci.\ USA} \textbf{102}, 13019 (2005).
\bibitem{Israelachvili2011} J.~N.~Israelachvili, \textit{Intermolecular and
Surface Forces}, 3rd ed.\ (Academic Press, 2011).
\bibitem{Weiss2008} U.~Weiss, \textit{Quantum Dissipative Systems},
3rd ed.\ (World Scientific, 2008).
\bibitem{Kubo1966} R.~Kubo,
\textit{Rep.\ Prog.\ Phys.} \textbf{29}, 255 (1966).
\bibitem{BreuerPetruccione2002} H.-P.~Breuer and F.~Petruccione,
\textit{The Theory of Open Quantum Systems} (Oxford Univ.\ Press, 2002).
\bibitem{Ishizaki2009} A.~Ishizaki and G.~R.~Fleming,
\textit{J.\ Chem.\ Phys.} \textbf{130}, 234111 (2009).
\bibitem{Tanimura2020} Y.~Tanimura,
\textit{J.\ Chem.\ Phys.} \textbf{153}, 020901 (2020).
\bibitem{Nogales1998} E.~Nogales, S.~G.~Wolf, and K.~H.~Downing,
\textit{Nature} \textbf{391}, 199 (1998).
\bibitem{Adolphs2006} J.~Adolphs and T.~Renger,
\textit{Biophys.\ J.} \textbf{91}, 2778 (2006).
\bibitem{Sept2003} D.~Sept, N.~A.~Baker, and J.~A.~McCammon,
\textit{Protein Sci.} \textbf{12}, 2257 (2003).
\bibitem{Moix2017} J.~M.~Moix \textit{et al.},
\textit{J.\ Phys.\ Chem.\ B} \textbf{121}, 3024 (2017).
\bibitem{Zurek2003} W.~H.~Zurek,
\textit{Rev.\ Mod.\ Phys.} \textbf{75}, 715 (2003).
\bibitem{Holstein1959} T.~Holstein,
\textit{Ann.\ Phys.\ (NY)} \textbf{8}, 325 (1959).
\bibitem{Mukamel1995} S.~Mukamel, \textit{Principles of Nonlinear Optical
Spectroscopy} (Oxford Univ.\ Press, 1995).
\bibitem{Hamm2011} P.~Hamm and M.~Zanni, \textit{Concepts and Methods of
2D Infrared Spectroscopy} (Cambridge Univ.\ Press, 2011).
\bibitem{Ollivier2001} H.~Ollivier and W.~H.~Zurek,
\textit{Phys.\ Rev.\ Lett.} \textbf{88}, 017901 (2001).
\bibitem{Luo2008} S.~Luo,
\textit{Phys.\ Rev.\ A} \textbf{77}, 042303 (2008).
\bibitem{Knill1998} E.~Knill and R.~Laflamme,
\textit{Phys.\ Rev.\ Lett.} \textbf{81}, 5672 (1998).
\bibitem{Datta2008} A.~Datta, A.~Shaji, and C.~M.~Caves,
\textit{Phys.\ Rev.\ Lett.} \textbf{100}, 050502 (2008).
\bibitem{Borgers2003} C.~B\"orgers and N.~Kopell,
\textit{Neural Comput.} \textbf{15}, 509 (2003).
\bibitem{Brunel2003} N.~Brunel and X.-J.~Wang,
\textit{J.\ Neurophysiol.} \textbf{90}, 415 (2003).
\bibitem{Buzsaki2012} G.~Buzs\'aki and X.-J.~Wang,
\textit{Annu.\ Rev.\ Neurosci.} \textbf{35}, 203 (2012).
\bibitem{Fries2001} P.~Fries \textit{et al.},
\textit{Science} \textbf{291}, 1560 (2001).
\bibitem{Buschman2007} T.~J.~Buschman and E.~K.~Miller,
\textit{Science} \textbf{315}, 1860 (2007).
\bibitem{Cardin2009} J.~A.~Cardin \textit{et al.},
\textit{Nature} \textbf{459}, 663 (2009).
\bibitem{Fries2015} P.~Fries,
\textit{Neuron} \textbf{88}, 220 (2015).
\bibitem{Poirazi2020} P.~Poirazi and A.~Papoutsi,
\textit{Nat.\ Rev.\ Neurosci.} \textbf{21}, 303 (2020).
\bibitem{Galarreta1999} M.~Galarreta and S.~Hestrin,
\textit{Nature} \textbf{402}, 72 (1999).
\bibitem{Bennett2004} M.~V.~L.~Bennett and R.~S.~Zukin,
\textit{Neuron} \textbf{41}, 495 (2004).
\bibitem{Anastassiou2011} C.~A.~Anastassiou \textit{et al.},
\textit{Nat.\ Neurosci.} \textbf{14}, 217 (2011).
\bibitem{Frohlich2010} F.~Fr\"ohlich and D.~A.~McCormick,
\textit{Neuron} \textbf{67}, 129 (2010).
\bibitem{Santelices2017} I.~B.~Santelices \textit{et al.},
\textit{Sci.\ Rep.} \textbf{7}, 9594 (2017).
\bibitem{Penrose1994} R.~Penrose, \textit{Shadows of the Mind} (Oxford
Univ.\ Press, 1994).
\bibitem{Hameroff2014} S.~Hameroff and R.~Penrose,
\textit{Phys.\ Life Rev.} \textbf{11}, 39 (2014).
\bibitem{Kerskens2022} C.~M.~Kerskens and D.~L\'opez~P\'erez,
\textit{J.\ Phys.\ Commun.} \textbf{6}, 105001 (2022).
\bibitem{Warren2023} W.~S.~Warren,
\textit{J.\ Phys.\ Commun.} \textbf{7}, 018001 (2023).
\bibitem{Borst1999} A.~Borst and F.~E.~Theunissen,
\textit{Nat.\ Neurosci.} \textbf{2}, 947 (1999).
\bibitem{Strong1998} S.~P.~Strong \textit{et al.},
\textit{Phys.\ Rev.\ Lett.} \textbf{80}, 197 (1998).
\bibitem{Giovannetti2004} V.~Giovannetti, S.~Lloyd, and L.~Maccone,
\textit{Science} \textbf{306}, 1330 (2004).
\bibitem{NeymanPearson1933}
J.~Neyman and E.~S.~Pearson,
\textit{On the Problem of the Most Efficient Tests of Statistical Hypotheses},
\textit{Phil.\ Trans.\ R.\ Soc.\ Lond.\ A} \textbf{231}, 289 (1933).

\bibitem{Cramer1946}
H.~Cramér,
\textit{Mathematical Methods of Statistics}
(Princeton University Press, Princeton, 1946).

\bibitem{Rao1945}
C.~R.~Rao,
\textit{Information and the Accuracy Attainable in the Estimation of
Statistical Parameters},
\textit{Bull.\ Calcutta Math.\ Soc.} \textbf{37}, 81 (1945).

\bibitem{Kay1993}
S.~M.~Kay,
\textit{Fundamentals of Statistical Signal Processing, Volume I:
Estimation Theory}
(Prentice Hall, 1993).

\bibitem{Schwarz1978}
G.~Schwarz,
\textit{Estimating the Dimension of a Model},
\textit{Ann.\ Stat.} \textbf{6}, 461 (1978).
\end{thebibliography}



\end{document}